\documentclass[11pt]{article}
\usepackage[a4paper,margin=25mm]{geometry}
\usepackage[T1]{fontenc}
\usepackage{amsmath,amssymb,amsthm,mathtools,mathrsfs,longtable,array,booktabs}
\usepackage{microtype}
\usepackage[hidelinks,hypertexnames=false]{hyperref}
\newtheorem{theorem}{Theorem}
\newtheorem{proposition}[theorem]{Proposition}
\newtheorem{lemma}[theorem]{Lemma}
\newtheorem{corollary}[theorem]{Corollary}
\theoremstyle{definition}
\newtheorem{definition}[theorem]{Definition}
\newtheorem{remark}[theorem]{Remark}
\providecommand{\R}{\mathbb R}
\providecommand{\C}{\mathbb C}
\providecommand{\Z}{\mathbb Z}
\providecommand{\N}{\mathbb N}
\providecommand{\dd}{\,d}
\title{Original zeta, shifted theta sums, and retained support}
\author{Programme derivations from the cited human sources}
\date{23 September 2026}
\begin{document}
\maketitle
The original functions, their complete transform multipliers, both
endpoint terms, boundary components, exceptional local lattices and
support labels are retained. The scalar kernel is recovered by an
explicit inverse Mellin integral. Neither this scalar inverse nor a
coefficient lift is asserted to reconstruct the entire arithmetic base.
\tableofcontents
\clearpage

% Complete source: ORIGINAL_ZETA_THETA_CORRECTION.tex
% ZTC: original zeta, original shifted arithmetic family, retained theta data.
\section{Original zeta and the retained theta comparison}
\label{sec:original-zeta-theta-correction}

\paragraph{Correction and provenance.}
The statement that the Rodgers--Tao heat kernel was the programme's
original construction over the entire absolute base was not established.
THS1--31 begin with that classical analytic kernel and construct a
support-preserving coefficient lift. HTB1--35 establish its exact
comparison with the shifted Hurwitz family. Neither construction is a
derivation of the scalar kernel from the entire arithmetic geometry.
The distinction is measurable: THS22 already proves that all supported
zero times act identically on the actual theta orbit. The ambient
coefficient injection does not make that orbit injective.

The original arithmetic functions in this calculation remain
\begin{equation}\tag{ZTC1}
 \zeta(s)=\sum_{n=1}^{\infty}n^{-s},\qquad
 F_t(s)=\sum_{n=0}^{\infty}(n+1+t)^{-s},
 \qquad \operatorname{Re}s>1,\quad t>-1.
\end{equation}
Their continuations, endpoints, and exact shift generator are proved in
HTB2--5, with the programme source identified there by version, hash,
and original theorem labels. In particular $F_0=\zeta$,
$F_1=\zeta-1$, and $\partial_tF_t(s)=-sF_t(s+1)$.
The latter family comes from the supplied April 2026 shifted-Hurwitz
programme draft, not from a proposed identification of two time variables.
Brad Rodgers and Terence Tao's
\href{https://arxiv.org/abs/1801.05914v5}{original author source},
equations \texttt{phidef}, \texttt{htdef}, \texttt{hoz}, and
\texttt{sas}, supplies the auxiliary heat kernel used in THS and HTB.
These are cited inputs, not new discoveries. The following comparisons
are derived explicitly. Their complete HTB, THS, and CFD proofs accompany
this chapter.

\subsection{A direct theta sum for the original arithmetic family}

For real $t>-1$ and $x>0$, retain the actual shifted frequencies and put
\begin{equation}\tag{ZTC2}
 K_t(x)=\sum_{n=0}^{\infty}
               e^{-\pi(n+1+t)^2x}.
\end{equation}
This is a scalar analytic representation of the stated arithmetic
family. No assertion that it determines the full support lattice is
part of this definition. For fixed $a=1+t>0$, monotonicity of the
summand as a function of $n$ gives
\begin{equation}\tag{ZTC3}
 K_t(x)\le e^{-\pi a^2x}
              +\int_0^\infty e^{-\pi(a+y)^2x}\,dy
 \le 1+\frac1{2\sqrt x}.
\end{equation}
For $x\ge1$ the same bound after retaining the first exponential
gives $K_t(x)\le C_a e^{-\pi a^2x}$, since
$(a+n)^2-a^2\ge n^2$ and
$C_a=\sum_{n\ge0}e^{-\pi((a+n)^2-a^2)}<\infty$.
Consequently the following integral converges absolutely for
$\operatorname{Re}s>1$:
\begin{equation}\tag{ZTC4}
 \int_0^\infty x^{s/2-1}K_t(x)\,dx
   =\pi^{-s/2}\Gamma(s/2)F_t(s),\qquad
 F_t(s)=\frac{\pi^{s/2}}{\Gamma(s/2)}
                \int_0^\infty x^{s/2-1}K_t(x)\,dx.
\end{equation}
Indeed, with $\sigma=\operatorname{Re}s>1$, the sum of the
absolute termwise integrals equals
$\pi^{-\sigma/2}\Gamma(\sigma/2)\sum_{n\ge0}(n+a)^{-\sigma}$.
Fubini therefore applies. Substitution $y=\pi(n+a)^2x$ in each
term proves the displayed identity with its full Gamma and pi factors.
Continuation outside that domain means the meromorphic continuation
proved in HTB4, not an assertion that the unsubtracted integral converges
there.

There is also an exact inverse recovering the kernel, not only the
arithmetic function. For every real $c>1$,
\begin{equation}\tag{ZTC4a}
 K_t(x)=\frac1{4\pi i}\int_{c-i\infty}^{c+i\infty}
       \pi^{-s/2}\Gamma(s/2)F_t(s)x^{-s/2}\,ds.
\end{equation}
Here the integral is absolutely convergent. To prove the inversion
with its constant, set $\sigma=c/2>0$ and
$g(v)=\exp(\sigma v-e^v)$. This function and all its derivatives
are integrable on $\mathbb R$: at $-\infty$ they are bounded
by a constant times $e^{\sigma v}$, and at $+\infty$ by a
polynomial in $e^v$ times $g(v)$. The substitution $u=e^v$
identifies its Fourier transform with $\Gamma(\sigma-i y)$.
Repeated integration by parts therefore bounds that Gamma factor
by $C_N(1+|y|)^{-N}$ for every fixed $N$.
Fourier inversion gives
$e^{-u}=(2\pi i)^{-1}\int_{\sigma-i\infty}^{\sigma+i\infty}
\Gamma(w)u^{-w}\,dw$. One may verify this inversion directly
by first multiplying the Fourier integrand by $e^{-\epsilon y^2}$:
the resulting expression is the convolution of $g$ with the
Gaussian approximate identity, which converges to $g$; the stated
integrable transform bound permits dominated convergence as
$\epsilon\downarrow0$. Apply the identity to
$u=\pi(n+1+t)^2x$, then put $s=2w$. The resulting factor is
$1/(4\pi i)$. Finally sum in $n$. Interchange is justified by
$\sum(n+1+t)^{-c}<\infty$ and the integrable Gamma bound.
This proves ZTC4a. ZTC4 and ZTC4a are mutual recovery formulas
for the specified original arithmetic family and its scalar
theta sums; they do not recover a support lattice from a scalar
function.

Differentiation in $t$ is locally uniform for $t$ in compact subsets
of $(-1,\infty)$ and $x$ in compact subsets of $(0,\infty)$, by
Gaussian majorants for a polynomial times $e^{-c n^2}$.
It gives the exact formula
\begin{equation}\tag{ZTC5}
 \partial_tK_t(x)=-2\pi x\sum_{n=0}^{\infty}
              (n+1+t)e^{-\pi(n+1+t)^2x}.
\end{equation}
For $\operatorname{Re}s>0$, integration of its absolute value is
bounded by a constant times $\sum(n+a)^{-\sigma-1}$, locally
uniformly in $t$. More precisely its Mellin transform is
\begin{equation}\tag{ZTC6}
 \int_0^\infty x^{s/2-1}\partial_tK_t(x)\,dx
 =-2\pi^{-s/2}\Gamma(s/2+1)F_t(s+1)
 =-s\pi^{-s/2}\Gamma(s/2)F_t(s+1).
\end{equation}
Thus the original arithmetic shift generator is represented directly,
without substituting the heat generator for it. The stated equality
extends meromorphically by HTB4--5.

The separated sheet and its entire missing summand are visible before
any transform:
\begin{equation}\tag{ZTC7}
 K_1(x)=K_0(x)-e^{-\pi x},\qquad
 D_t(s)=\zeta(s)-F_t(s),\qquad D_1(s)=1,
\end{equation}
\begin{equation}\tag{ZTC8}
 \pi^{-s/2}\Gamma(s/2)D_t(s)
   =\int_0^\infty x^{s/2-1}\bigl(K_0(x)-K_t(x)\bigr)\,dx
       \quad(\operatorname{Re}s>1).
\end{equation}
ZTC7 follows by deleting the $n=1$ summand of $K_0$; ZTC8 follows
by subtracting two absolutely convergent integrals. At $t=1$ the
integral is exactly $\pi^{-s/2}\Gamma(s/2)$, so its inverse returns
the actual constant $1$, including its amplitude and support when lifted.

\subsection{Original zeta with both Mellin endpoints retained}

At $t=0$, HTB6 proves
\begin{equation}\tag{ZTC9}
 K_0(x)=x^{-1/2}K_0(1/x)+\tfrac12x^{-1/2}-\tfrac12.
\end{equation}
Split ZTC4 at $x=1$, apply this identity only to the integral over
$(0,1)$, and substitute $y=1/x$ there. Initially
$\operatorname{Re}s>1$, and direct integration gives
\begin{equation}\tag{ZTC10}
 \zeta(s)=\frac{\pi^{s/2}}{\Gamma(s/2)}
 \left\{
 \frac1{s-1}-\frac1s
 +\int_1^\infty K_0(x)
             \left(x^{s/2-1}+x^{-(s+1)/2}\right)\,dx
 \right\}.
\end{equation}
The two separate rational terms are the transforms of the two separate
terms $x^{-1/2}/2$ and $-1/2$. The integral on $[1,\infty)$ is
entire in $s$, since exponential decay bounds every logarithmic
derivative on compact subsets of the $s$-plane. This proves the
meromorphic continuation of the identity and keeps both endpoint
contributions. Neither has been declared zero or absorbed into an
unrecorded constant.

For the original cosine heat kernel, retain
\begin{equation}\tag{ZTC11}
\begin{gathered}
 A(s)=\frac{s(s-1)}2\pi^{-s/2}\Gamma(s/2),\qquad
 s(z)=\frac12+\frac{iz}{2},\qquad z(s)=-2i(s-1/2),\\
 H_0(z)=\frac{s(z)(s(z)-1)}{16}
           \pi^{-s(z)/2}\Gamma(s(z)/2)\zeta(s(z)),\\
 \zeta(s)=\frac{16\pi^{s/2}}
                  {s(s-1)\Gamma(s/2)}H_0(z(s)).
\end{gathered}
\end{equation}
The proof is HTB6--10: the exact Mellin substitution and two
integrations by parts yield the displayed multiplier, with all
endpoint limits justified in the original convergence half-plane.
The last equality is an identity of meromorphic functions. It is
not a licence to replace the original zeta calculation by an entire
function calculation while dropping its multiplier.

\subsection{The boundary is part of the receiving object}

Let $\mathcal M_s$ and $\mathcal M_z$ be the complex vector spaces
of meromorphic functions on the respective complex planes. Retain
the linear isomorphism, already proved in HTB11,
\begin{equation}\tag{ZTC12}
 Uf(z)=\frac{A(s(z))}{8}f(s(z)),\qquad
 U^{-1}G(s)=\frac{8G(z(s))}{A(s)}.
\end{equation}
Define the joint map on the entire two-component space by
\begin{equation}\tag{ZTC13}
 \mathcal V:\mathcal M_s\oplus\mathcal M_s
          \longrightarrow\mathcal M_z\oplus\mathcal M_z,
 \qquad (F,D)\longmapsto\bigl(U(F+D),UD\bigr).
\end{equation}
It has the exact inverse
\begin{equation}\tag{ZTC14}
 \mathcal V^{-1}(H,J)
    =\left(\frac{8(H-J)\circ z}{A},\frac{8J\circ z}{A}\right).
\end{equation}
Substitution proves both compositions equal the identity; all entries
are meromorphic functions, so no undefined product of separate
singular point values is involved. On the actual arithmetic family,
\begin{equation}\tag{ZTC15}
\begin{gathered}
 J_t=UD_t,\qquad
 \mathcal V(F_t,D_t)=(H_0,J_t),\\
 F_t(s)=\frac{16\pi^{s/2}}
              {s(s-1)\Gamma(s/2)}
              \bigl(H_0(z(s))-J_t(z(s))\bigr),\\
 J_1(z)=\frac{A(s(z))}{8},\qquad
 J_t(i)=-\frac t8.
\end{gathered}
\end{equation}
For the last equality use HTB5:
$D_t(0)=\zeta(0)-F_t(0)=t$, $A(0)=-1$, and $z(0)=i$.
In particular the retained pair determines the real parameter $t$
by $t=-8J_t(i)$. All $t$ give the same first component $H_0$.
Discarding $J_t$ therefore identifies different actual arithmetic
inputs. This is an explicit failure of injectivity on the actual
programme family, not merely an example in an unrelated vector space.
On the unrestricted space the summation map $(G,J)\mapsto G+J$
has exactly the kernel $\{(B,-B):B\in\mathcal M_z\}$.

For the complete support carrier, use the common label on the
direct sum:
\begin{equation}\tag{ZTC16}
 G_L(\mathcal M_s\oplus\mathcal M_s)
   \xrightarrow{\mathcal V^\uparrow}
 G_L(\mathcal M_z\oplus\mathcal M_z),\qquad
 ((F,D),\lambda)\longmapsto(\mathcal V(F,D),\lambda).
\end{equation}
Here $G_L(V)=\{(0,\lambda):\lambda\in L\}\cup(V\times\{1_L\})$
with its original joins and meets. Additivity, compatibility with
the $G_L(\mathbb C)$ action, carrier membership, and the inverse
follow componentwise from linearity and ZTC14, as explicitly proved
in CFD11--12. This is a bijection for every bounded distributive $L$.
It does not silently replace the common label by independent labels
on two product factors. In particular $e=(0,1_L)$ and
$\tau=(0,0_L)$ remain distinct when $L$ is nontrivial.

\paragraph{The exact multiplication comparison.}
The preceding maps were declared as linear and semimodule maps.
Their injectivity does not identify the original multiplication with
ordinary multiplication of their images. This can be checked exactly,
without discarding the completion multiplier. Set
\begin{equation}\tag{ZTC16a}
 B(z)=\frac{A(s(z))}{8}
 =\frac{s(z)(s(z)-1)}{16}
       \pi^{-s(z)/2}\Gamma(s(z)/2).
\end{equation}
For meromorphic $f,g$ on the $s$-plane the two products are
\begin{equation}\tag{ZTC16b}
\begin{gathered}
 U(fg)(z)=B(z)f(s(z))g(s(z)),\\
 (Uf)(z)(Ug)(z)=B(z)^2f(s(z))g(s(z)).
\end{gathered}
\end{equation}
In particular $U1=B$, not the ordinary constant unit. The factor
$B$ is nonconstant: it has a zero at $z(1)=-i$ and poles at
$z(-2m)$, $m\ge1$. Thus ordinary multiplication is not preserved.
The exact transported product and its unit are
\begin{equation}\tag{ZTC16c}
 G\odot H=\frac{8G(z)H(z)}{A(s(z))},\qquad
 1_{\odot}(z)=\frac{A(s(z))}{8}.
\end{equation}
This formula is an identity of meromorphic functions, including
their germs at the zeros and poles of the multiplier. It is not
pointwise division of undefined singular values. Substitution
gives $(Uf)\odot(Ug)=U(fg)$. Directly,
$(G\odot H)\odot K=GHK/B^2=G\odot(H\odot K)$,
$B\odot G=G$, and distributivity follows from distributivity of
meromorphic multiplication. Together with ZTC12 this proves a
unital algebra isomorphism with the stated product.

For each local holomorphic ring $\mathcal O_{s_0}$ the affine
coordinate map identifies it with $\mathcal O_{z(s_0)}$.
Its exact image under $U$ is the embedded fractional lattice
$B\mathcal O_{z(s_0)}$. This lattice is closed under $\odot$,
since $(Bu)\odot(Bv)=Buv$, and has unit $B$.
Keeping instead the unmodified target $\mathcal O_{z(s_0)}$
can fail closure: at $s_0=1$, $B$ has a simple zero, so
$1\odot1=1/B$ is not holomorphic. This image-lattice statement
differs from the preimage lattice $A^{-1}\mathcal O$ in ZTC17:
the former transports an original holomorphic ring, whereas the
latter specifies which original meromorphic sections have a
holomorphic completed image.

The full support carrier admits the same exact comparison. Retain
its original addition, and define the target product by
\begin{equation}\tag{ZTC16d}
 (G,\lambda)\odot_L(H,\mu)
   =\left(\frac{8GH}{A\circ s},\lambda\wedge\mu\right),
 \qquad 1_{\odot_L}=(B,1_L).
\end{equation}
If the amplitude of a product is nonzero, both input amplitudes
are nonzero, hence both labels are $1_L$; this proves carrier
closure. If either amplitude is zero its output amplitude is
zero and the displayed meet retains its original label.
Associativity and distributivity follow from ZTC16c and the
meet/join laws of the bounded distributive lattice. The map
$(f,\lambda)\mapsto(Uf,\lambda)$ and its inverse therefore
give a unital semiring isomorphism from the original coefficient
semiring to this specified target, preserving $e$ and $\tau$.
An ideal $P$ is sent to its literal image $U^{\uparrow}(P)$;
closure under addition and multiplication follows by applying the
isomorphism. Its inverse proves the converse. Moreover a product
belongs to that image precisely when the corresponding original
product belongs to $P$, proving that prime ideals correspond.
This proves the exact spectrum comparison for these specified
coefficient semirings. It is not a derivation of the entire
arithmetic base from the scalar function $H_0$.

For ZTC13--16, no multiplication on pairs was specified: the
proved assertion there remains a common-label semimodule
isomorphism. Neither ordinary multiplication nor a Weil pairing
is silently inferred from that assertion.

\subsection{Local zeros, poles, and the original trace}

The earlier conversational assertion that the singular module map
had not been established was also too broad: HTB13 already proves
$A^{-1}\mathcal O\xrightarrow{U}\mathcal O$ with the affine
coordinate change. CFD1--10 now make the comparison with the
\emph{original} regular lattice and every finite jet explicit.
At a finite point let $h=s-s_0$ and $A=h^k a(h)$, $a(0)\ne0$.
Retain $I=A^{-1}\mathcal O=h^{-k}\mathcal O$ as an embedded
submodule of the meromorphic germs, and $J=I\cap\mathcal O$.
For nonzero $f\in I$ the original order is
\begin{equation}\tag{ZTC17}
 \operatorname{ord}_{s_0}f
 =\dim_{\mathbb C}(I/\mathcal O f)
  -\dim_{\mathbb C}(I/J)
  +\dim_{\mathbb C}(\mathcal O/J).
\end{equation}
This is proved by the power-of-$h$ bases in CFD6--7, without
identifying the reference lattices. At a trivial zero $s_0=-2m$,
$I=h\mathcal O$, $I/\mathcal O\zeta=0$ and
$\mathcal O/I$ has dimension $1$. At $s_0=1$,
$I=h^{-1}\mathcal O$, $I/\mathcal O\zeta=0$ and
$I/\mathcal O$ has dimension $1$. Their signs in ZTC17 give,
respectively, $+1$ and $-1$. At $s_0=0$, $A(0)=-1$ is a unit.
CFD4 provides the invertible triangular matrices recovering all
Laurent coefficients with the exact factors $8$ and $i/2$.
Keeping only the zero algebra of the completed section would give
zero in both cancelled cases; keeping these comparison modules
restores the actual divisor.

The earlier conversational explanation of the fibres requires a
further correction. The map
\begin{equation}\tag{ZTC17a}
 \mathcal O/h\mathcal O\longrightarrow
 h\mathcal O/h^2\mathcal O,\qquad [u]\longmapsto[hu]
\end{equation}
is an isomorphism of $\mathcal O$-modules. It is well-defined because
$h(h\mathcal O)=h^2\mathcal O$, surjective by the definition of
$h\mathcal O$, and injective because $hu\in h^2\mathcal O$ implies
$u\in h\mathcal O$ in the local integral domain. Its inverse is
$[hu]\mapsto[u]$. Consequently the two abstract fibres alone do
not account for the loss. What must be retained is their specified
embedding and the original reference lattice: multiplication by
$A=h^{-1}a(h)$ identifies $I=h\mathcal O$ with $\mathcal O$, but
the length-one quotient $\mathcal O/I$ records the cancelled
original zero. Replacing the entire marked comparison by the
completed section's zero quotient discards precisely that datum.
Likewise multiplication by nonzero meromorphic $A$ is an
automorphism of the meromorphic sheaf as a module. It is not
a ring automorphism for ordinary multiplication, as ZTC16b proves.
It is not an automorphism
of the fixed holomorphic lattice at its zeros and poles. These
category statements do not contradict one another.

Logarithmic differentiation gives the original trace identity
\begin{equation}\tag{ZTC18}
 \frac{\zeta'(s)}{\zeta(s)}
 =-2i\frac{H_0'(z(s))}{H_0(z(s))}
 -\frac1s-\frac1{s-1}+\frac12\log\pi-\frac12\psi(s/2).
\end{equation}
It first holds away from zeros and poles and then meromorphically;
$\psi=\Gamma'/\Gamma$. For a positively oriented finite contour
$\gamma$ avoiding all zeros and poles of the functions and all
singularities of every displayed integrand, and a holomorphic test $g$
on and inside it, multiplication by $g$ and integration gives
\begin{equation}\tag{ZTC19}
\begin{split}
 \frac1{2\pi i}\int_\gamma g\,\frac{\zeta'}\zeta\,ds
 ={}&\frac1{2\pi i}\int_\gamma
           g(s)(-2i)\frac{H_0'(z(s))}{H_0(z(s))}\,ds\\
 &-\frac1{2\pi i}\int_\gamma g(s)
   \left(\frac1s+\frac1{s-1}-\frac12\log\pi
                         +\frac12\psi(s/2)\right)ds.
\end{split}
\end{equation}
At every zero or pole, write a meromorphic function as
$(s-s_0)^r$ times a holomorphic unit. Its logarithmic derivative
then has residue $r$. The residue theorem proves that ZTC19
counts the original zeros with multiplicity and the original poles
with negative multiplicity, including every trivial zero inside
the contour. No infinite-contour limit or decay assumption is
hidden here. Extending this finite identity to a particular infinite
trace requires the actual convergence estimates of that trace.

\subsection{Two different trajectories through the same original zeta}

The original-$s$ function associated with the classical heat path is
\begin{equation}\tag{ZTC20}
 f_h^{\rm heat}(s)=\frac{16\pi^{s/2}}
                         {s(s-1)\Gamma(s/2)}H_h(z(s)),
 \qquad f_0^{\rm heat}=\zeta.
\end{equation}
CFD18--21 give its full receiving operator and prove
\begin{equation}\tag{ZTC21}
 \left.\partial_hf_h^{\rm heat}(-2)\right|_{h=0}=0,
 \qquad
 \left.\partial_tF_t(-2)\right|_{t=0}=-\frac16.
\end{equation}
For completeness, the heat equation and chain rule give
$\partial_hf=(Af)''/(4A)$. At $h=0$, $A\zeta$ is holomorphic
at $-2$ whereas $1/A$ has a simple zero there, giving the first
value. HTB5 gives the second as $2\zeta(-1)=-1/6$.
The discrepancy holds on the actual initial arithmetic function.
CFD21 further proves that all the zeros at $-2m$ remain simple
along the real heat path, while its residue at $1$ varies strictly;
the actual Hurwitz path has constant residue $1$ and values
$-B_{2m+1}(1+t)/(2m+1)$ at those points.

Thus the valid original-$\zeta$ receiver consists of the original
functions, their exact transform and inverse, the retained boundary
component, the embedded local comparison modules, and all support
labels. The scalar theta orbit by itself is a smaller observable.
These proofs neither exhibit an off-line zero of the original
$\zeta$ nor rule one out. They establish precisely which input and
which deformation a subsequent estimate is estimating.

\clearpage

% Complete source: COMPLETION_FAITHFULNESS_DERIVATION.tex
% Independent derivation, 2026-09-23. Prefix CFD.
% Sources actually read: HTB in full; THS in full; HZA in full.
\section{Completion, local lattices, and the exact observable of support}
\label{sec:completion-faithfulness-independent}

\paragraph{Sources and scope.}
The analytic source is Brad Rodgers and Terence Tao,
\href{https://arxiv.org/abs/1801.05914v5}{\emph{The de Bruijn--Newman constant is non-negative}},
original author source equations \texttt{phidef}, \texttt{htdef},
\texttt{hoz}, and \texttt{sas}. Its kernel, coordinates, and Mellin
identity are proved directly in the included HTB1--10; the convergence
and positivity needed here are proved in THS1--15. The shifted Hurwitz
family is the supplied April 2026 programme draft, byline ``Compiled
theorem draft'', equations \texttt{thm:universal-flow} and
\texttt{thm:zero-motion}, with no inferred human attribution.
The full support carrier and its public source are stated in HTB29--31
and HZA1. The complete HTB, THS, and HZA sources were read for this
derivation. All new comparisons below are proved here. A lift of a
classical analytic function to that carrier is not thereby a geometric
construction of that function from the entire arithmetic base.

\subsection{The fractional module and every finite jet}

Retain the exact constants
\begin{equation}\tag{CFD1}
 A(s)=\frac{s(s-1)}2\pi^{-s/2}\Gamma(s/2),\quad
 s(z)=\frac12+\frac{iz}{2},\quad z(s)=-2i(s-1/2),\quad
 (Uf)(z)=\frac{A(s(z))}{8}f(s(z)).
\end{equation}
Fix a finite $s_0$, set $h=s-s_0$, $z_0=z(s_0)$, and $w=z-z_0$.
Write $A=h^k a(h)$, where $k\in\mathbb Z$ and $a$ is a holomorphic
unit. Let $\mathcal O_s$ and $\mathcal O_z$ be the respective rings of
holomorphic germs, and $\mathcal M_s$ the field of meromorphic germs.
Define the specified embedded fractional module
\begin{equation}\tag{CFD2}
 I=A^{-1}\mathcal O_s=h^{-k}\mathcal O_s\ \subset\mathcal M_s.
 \qquad U:I\longrightarrow\mathcal O_z,
 \qquad U^{-1}g=8g(z(s))/A(s).
\end{equation}
The equality of fractional modules uses that $a$ is a unit. The two
displayed substitutions compose to the identity. Thus this is an
isomorphism of modules over the coordinate isomorphism
$C:\mathcal O_s\to\mathcal O_z$, $C(h)=iw/2$. In particular
\begin{equation}\tag{CFD3}
 U(h^{N+1}I)=w^{N+1}\mathcal O_z,\qquad
 I/h^{N+1}I\xrightarrow{\ \overline U\ }
                  \mathcal O_z/w^{N+1}\mathcal O_z
 \quad(N\ge0)
\end{equation}
are isomorphisms. No undefined product of point values is used at a
pole or zero of $A$.

For the exact coefficient matrices, write
$a(h)=\sum_{j\ge0}a_jh^j$ and
$f=h^{-k}\sum_{j\ge0}f_jh^j$. The coefficient of $w^n$ in $Uf$ is
\begin{equation}\tag{CFD4}
 g_n=\frac18\left(\frac i2\right)^n
            \sum_{j=0}^n a_j f_{n-j},\qquad
 f_n=\sum_{j=0}^n d_j(2/i)^{n-j}g_{n-j},\qquad
 \sum_{j\ge0}d_jh^j=\frac8{a(h)}.
\end{equation}
The first formula is multiplication of the two convergent power
series followed by $h=iw/2$; the second reverses those steps. On
jets through degree $N$ the matrix is triangular with diagonal
$a_0(i/2)^n/8$, which never vanishes. These are the original coordinate
and amplitude factors, not a replacement of $h$ or $w$ by a rescaled
coordinate.

The associated graded map in degree $n$ sends the class of
$h^{n-k}$ to $a_0(i/2)^n w^n/8$. Consequently the completed order is
\begin{equation}\tag{CFD5}
 \operatorname{ord}_{z_0}(Uf)=k+\operatorname{ord}_{s_0}(f).
\end{equation}
The germ $f$ and all of its Laurent coefficients are recovered from
$Uf$ once the embedded module $I$, the specified multiplier $A$, and
the coordinate map are retained.

\subsection{Which quotient remembers a cancelled divisor}

For every nonzero $f\in I$ there is an exact module isomorphism
\begin{equation}\tag{CFD6}
 I/\mathcal O_s f\ \xrightarrow{\ \overline U\ }\
 \mathcal O_z/\mathcal O_z(Uf).
 \qquad
 \dim_{\mathbb C}(I/\mathcal O_s f)
       =k+\operatorname{ord}_{s_0}f.
\end{equation}
Indeed $U(qf)=C(q)Uf$ and $C$ is surjective. The right hand quotient
is the quotient by a unit times $w^{k+\operatorname{ord}f}$, whose
classes $1,w,\ldots,w^{k+\operatorname{ord}f-1}$ form a basis.
This quotient measures the zero of the completed section in $I$.
It does not, by itself, measure the zero or pole of $f$ relative to
the different reference lattice $\mathcal O_s\subset\mathcal M_s$.

The exact comparison retains both embeddings:
\begin{equation}\tag{CFD7}
\begin{gathered}
 J=I\cap\mathcal O_s=h^{\max(0,-k)}\mathcal O_s,\qquad
 I+\mathcal O_s=h^{\min(0,-k)}\mathcal O_s,\\
 \dim_{\mathbb C}(I/J)=\max(k,0),\qquad
 \dim_{\mathbb C}(\mathcal O_s/J)=\max(-k,0),\\
 \operatorname{ord}(f)
 =\dim(I/\mathcal O_s f)-\dim(I/J)+\dim(\mathcal O_s/J).
\end{gathered}
\end{equation}
The first line follows by comparing integer orders of germs; the
second by bases of the relevant powers of $h$; the third follows
from CFD6 and $\max(k,0)-\max(-k,0)=k$. Thus this object retains
the cancelled divisor with its sign and its precise local module.

For the actual multiplier, HTB18 gives
\begin{equation}\tag{CFD8}
\begin{gathered}
 A(0)=-1,\quad A(s)=\tfrac12(s-1)+O((s-1)^2),\\
 \operatorname*{Res}_{s=-2m}A(s)
    =a_m=\frac{2(-1)^m(2m+1)\pi^m}{(m-1)!}\ne0
       \quad(m\ge1).
\end{gathered}
\end{equation}
There are no other finite zeros or poles. At $s_0=-2m$, let
$h=s+2m$. Then $k=-1$, $I=h\mathcal O_s$, and the Riemann zeta
function is $\zeta=h c(h)$ with $c(0)\ne0$. To see the last
nonvanishing without an additional zero-location assertion, HTB9
gives $\xi(-2m)=\xi(1+2m)$, and the latter is nonzero because all
factors of $A(1+2m)\zeta(1+2m)$ are positive. Hence
\begin{equation}\tag{CFD9}
 I/\mathcal O_s\zeta=0,\quad
 \mathcal O_s/I\cong\mathcal O_s/(h),\quad
 (U\zeta)(z_0)=a_m\zeta'(-2m)/8.
\end{equation}
The vanished quotient in the first expression cannot serve as the
record of the cancelled simple zero. The second expression does
record its length, and CFD4 recovers every derivative from the
completed germ. The one-dimensional module $h\mathcal O_s/h^2
\mathcal O_s$ is the degree-zero associated graded of $I$; it
contains the leading zero coefficient and maps isomorphically to
$\mathcal O_z/w\mathcal O_z$. It is not the same quotient as
$I/\mathcal O_s\zeta$.

At $s_0=1$, $k=1$, $I=h^{-1}\mathcal O_s$ and
$\zeta=h^{-1}(1+O(h))$, so
\begin{equation}\tag{CFD10}
 I/\mathcal O_s\zeta=0,\quad
 I/\mathcal O_s\cong h^{-1}\mathcal O_s/\mathcal O_s,
 \quad (U\zeta)(-i)=1/16.
\end{equation}
Here the one-dimensional comparison quotient records the cancelled
pole. At $s=0$ the multiplier is a unit and no divisor cancellation
occurs. At all points of the open critical strip it is a unit,
so the ordinary local zero algebras of $\zeta$ and $U\zeta$ are
isomorphic by coordinate substitution, with every multiplicity
unchanged. CFD9--10 explain precisely why this strip assertion
does not replace the global comparison.

\subsection{All support labels on modules and quotients}

For each complex vector space $V$, use exactly
\begin{equation}\tag{CFD11}
 G_L(V)=\{(0,\lambda):\lambda\in L\}\cup(V\times\{1_L\}),
 \quad (v,\lambda)+(v',\mu)=(v+v',\lambda\vee\mu),
 \quad(c,\alpha)(v,\lambda)=(cv,\alpha\wedge\lambda).
\end{equation}
Every complex-linear map $T$ above lifts by
$T^\uparrow(v,\lambda)=(Tv,\lambda)$. The carrier condition holds
because $T0=0$; addition and scalar action commute componentwise.
The inverse lifts when $T$ is invertible. This proves all module
and finite-jet isomorphisms above over the entire bounded
distributive lattice $L$, with no finiteness or Boolean restriction.

For a subspace $W\subset V$, the quotient map has the precise
congruence
\begin{equation}\tag{CFD12}
 (v,\lambda)\sim(v',\lambda')
 \ \Longleftrightarrow\
 \lambda=\lambda'\ \hbox{ and }\ v-v'\in W,
 \qquad
 G_L(V)/{\sim}\ \cong G_L(V/W).
\end{equation}
The relation is preserved by addition and scalar action, since $W$
is a subspace and the same labels undergo the same joins or meets.
Its classes are exactly the fibres of $(v,\lambda)\mapsto
(v+W,\lambda)$, which is surjective: a nonzero coset has top label,
and each labelled zero is the image of the corresponding labelled
zero. This proves the quotient assertion. In particular
$G_L(0)=\{(0,\lambda):\lambda\in L\}$ still retains every label;
it contains no amplitude dimension recording a cancelled divisor.
The defects in CFD7--10 must therefore remain even after this lift.

With the original $K_L=K\otimes_{\mathbb B}L$, the map
\begin{equation}\tag{CFD13}
 \Phi_V(v,\lambda)=(v,1_K\otimes\lambda)
 \quad\hbox{satisfies}\quad
 (T\times\mathrm{id})\Phi_V=\Phi_W T^\uparrow.
\end{equation}
The zero/nonzero character of the tropical semifield $K$ tensored
with $L$ retracts $1_K\otimes\lambda$ to $\lambda$, proving
injectivity. Thus neither completion nor these module comparisons
identify $e=(0,1_L)$ with $\tau=(0,0_L)$. Multiplication by $A/8$
is a module map, not an ordinary unital algebra map. The ordinary
product transports to $(G\odot H)(z)=8G(z)H(z)/A(s(z))$ and unit
$A(s(z))/8$, as in HTB12. Statements about prime spectra require
that declared transported product, rather than the ordinary
pointwise product.

\subsection{The theta orbit and a faithful supported-time probe}

Keep the actual coefficient vector of THS18:
$a_{2j}(t)=(h_j(t),1_L)$ and $a_{2j+1}(t)=\tau$.
For $x=(r,\lambda)\in G_L(\mathbb C)$ the exact matrix is
\begin{equation}\tag{CFD14}
 U_x(i,j)=
 \begin{cases}
 (1,1_L),&j=i,\\
 ((-r)^m j!/(m!i!),\lambda),&j=i+2m,\ m\ge1,\\
 \tau,&\text{otherwise}.
 \end{cases}
\end{equation}
THS10--12 give the absolute amplitude convergence. The diagonal
summand has top support at every even index, while every odd
summand is $\tau$. Hence
\begin{equation}\tag{CFD15}
 U_{(r,\lambda)}\mathbf a(t)=\mathbf a(t+r),\qquad
 U_{(0,\lambda)}\mathbf a(t)=\mathbf a(t)
                    \quad\hbox{for every }\lambda\in L.
\end{equation}
For nontrivial $L$ this orbit observable is not injective in the
supported-time parameter. Keeping the label in the ambient carrier
does not make this particular observable distinguish it.

The full matrix representation nevertheless is injective:
\begin{equation}\tag{CFD16}
 U_x(0,2)=(-2r,\lambda),\qquad
 x=(-\tfrac12,1_L)\,U_x(0,2).
\end{equation}
Both the amplitude and every allowed label are recovered by this
entry. Equivalently, let $p$ be the supported polynomial whose
only non-$\tau$ coefficient is $p_2=(1,1_L)$. Direct multiplication
gives
\begin{equation}\tag{CFD17}
 (U_xp)_0=(-2r,\lambda),\quad (U_xp)_2=(1,1_L),\quad
 (U_xp)_j=\tau\ (j\ne0,2).
\end{equation}
Thus the degree-two probe restores supported-time faithfulness;
its left inverse is the constant-coefficient map followed by
multiplication by $(-1/2,1_L)$. Its sums are finite. A degree-zero
or degree-one polynomial cannot provide this test because the
second derivative annihilates both and the strictly upper even
entries have no source coefficient. The proof establishes the
parameter observable just specified, not faithfulness of a
construction of the complete arithmetic geometry.

Another proved loss of information is the finite-cardinality map
of HZA2: for finite $L$, $|G_L(\mathbb Z/n\mathbb Z)|=n-1+|L|$
and its Dirichlet series is $\zeta(s,|L|)$. A four-element chain
and the four-element Boolean lattice give the same series but
different meet/join structures. Hence this scalar receiver is
also not a reconstruction of the full lattice. The injective
$\Phi_V$ and the count are different, explicitly specified maps.

\subsection{Completion does not identify the two deformations}

On meromorphic functions let $Df(s)=s f(s+1)$ and
$B=-\partial_z^2$. The pullback of the heat generator is
\begin{equation}\tag{CFD18}
 \mathcal L_H=U^{-1}BU,\qquad
 \mathcal L_H f=\frac1{4A}(Af)'',\qquad
 (\mathcal L_H+D)f=
 \frac14 f''+\frac{A'}{2A}f'
             +\frac{A''}{4A}f+s f(s+1).
\end{equation}
This follows by the chain rule $ds/dz=i/2$ and differentiating
twice; the factors $8$ cancel. The Hurwitz generator is $-D$.
On the constant function $1$, their difference is
$A''/(4A)+s$. At each $s_0=-2m$, writing $A=h^{-1}a(h)$ gives
$A''/A=2h^{-2}+O(h^{-1})$, so the difference has a nonzero
double-pole coefficient $1/2$. It is not the zero operator.

It also differs on the actual initial function $\zeta$. Near
$s=-2$, the residue calculation HTB18 gives
$A(s)=-6\pi/(s+2)+O(1)$. Since $A\zeta=\xi$ is entire,
\begin{equation}\tag{CFD19}
 \mathcal L_H\zeta(-2)=0,\qquad
 (-D\zeta)(-2)=2\zeta(-1)=-\tfrac16,\qquad
 ((\mathcal L_H+D)\zeta)(-2)=\tfrac16.
\end{equation}
Here $\zeta(-1)=-1/12$ follows directly from HTB5 and
$B_2(1)=1/6$. Therefore the two tangent vectors are different
already before completion. Their transported difference has
residue $-\pi/8$ in the $s$ coordinate for $U$ without coordinate
substitution, and $-\pi/(4i)$ in the $z$ coordinate. Equivalently
the Hurwitz tangent alone has the opposite residue $\pi/(4i)$,
whereas the heat tangent is entire, as in HTB21.

This distinction persists on the exact trajectories. Write
\begin{equation}\tag{CFD20}
 f_h^{\mathrm{heat}}(s)=\frac{8H_h^{\mathrm{heat}}(z(s))}{A(s)}.
\end{equation}
For every real $h$ and every $c\in\mathbb R$,
$H_h^{\mathrm{heat}}(ic)=\int_0^\infty
e^{hu^2}\Phi(u)\cosh(cu)\,du>0$; convergence and strict
positivity follow from THS3--6. It follows that
\begin{equation}\tag{CFD21}
\begin{gathered}
 \operatorname{ord}_{s=-2m}f_h^{\mathrm{heat}}=1\quad(m\ge1),
 \qquad
 \operatorname*{Res}_{s=1}f_h^{\mathrm{heat}}
                       =16H_h^{\mathrm{heat}}(-i),\\
 \frac{d}{dh}\operatorname*{Res}_{s=1}f_h^{\mathrm{heat}}
       =16\int_0^\infty u^2e^{hu^2}\Phi(u)\cosh u\,du>0.
\end{gathered}
\end{equation}
The first conclusion uses the simple zero of $1/A$ and the
strictly positive numerator; the residue uses $A'(1)=1/2$;
the derivative is justified by THS6. At $h=0$ the residue is
$1$. In contrast $F_t(s)=\zeta(s,1+t)$ has residue $1$ at
$s=1$ for every allowed $t$ and values
$F_t(-2m)=-B_{2m+1}(1+t)/(2m+1)$. Already its first completed
pole has residue $2\pi(1+t)(t+1/2)t$. Thus completion is
invertible on the specified meromorphic data, but the heat
trajectory preserves the classical trivial-zero locations by
its construction and is not the arithmetic Hurwitz trajectory.

The exact bridge remains HTB26--28:
$F_t+(\zeta-F_t)=\zeta$ and
$UF_t+U(\zeta-F_t)=U\zeta$. It retains the complete principal
parts, reflection defect, and every support label. Omitting the
second summand would change the arithmetic input; retaining it
is an identity returning the initial function, not a proof that
the two distinct generators agree or that either trajectory
reconstructs the full arithmetic base.

\subsection{Faithful parameter representation is not invertible heat transport}

The injectivity of $x\mapsto U_x$ proved in CFD16 does not mean
that each $U_x$ is an injective endomorphism of supported polynomials.
In fact, for every $x=(r,\lambda)\ne\tau$ one has $\lambda\ne0_L$.
Take the degree-two probe $p$ of CFD17 and let $p'$ have the same
coefficients except $p'_0=(0,\lambda)$ instead of $p_0=\tau$.
They are distinct. Their degree-zero images are respectively
$(-2r,\lambda)$ and
$(0,\lambda)+(-2r,\lambda)=(-2r,\lambda)$, and every other image
coefficient agrees. Thus
\begin{equation}\tag{CFD22}
 p\ne p',\qquad U_xp=U_xp'\quad(x\ne\tau).
\end{equation}
So no such $U_x$ has an inverse on the full supported polynomial
space. The classical time-group relation on the actual saturated
theta vector results from CFD15; the full supported-time family
instead satisfies the additive monoid relation of THS20. This is
an exact support effect, despite injectivity of its parameter map.

\subsection{The direct shifted theta kernel before completion}

For real $t>-1$, put $a=1+t>0$ and retain the direct kernel
\begin{equation}\tag{CFD23}
 K_t(x)=\sum_{n=0}^{\infty}e^{-\pi(n+a)^2x},\quad x>0,
 \qquad
 K_1(x)=K_0(x)-e^{-\pi x}.
\end{equation}
The series and its parameter derivatives converge locally uniformly
for $x>0$ and $a$ in bounded real intervals, by the Gaussian decay in
$n$. The last identity deletes precisely the $n=0$ term of the
$a=1$ series. For $\operatorname{Re}s>1$ one has
\begin{equation}\tag{CFD24}
 \int_0^\infty K_t(x)x^{s/2-1}\,dx
       =\pi^{-s/2}\Gamma(s/2)\zeta(s,a).
\end{equation}
Indeed, if $\sigma=\operatorname{Re}s>1$, the sum of integrals of
the absolute values is
$\pi^{-\sigma/2}\Gamma(\sigma/2)
\sum_{n\ge0}(n+a)^{-\sigma}<\infty$.
Fubini's theorem applies. Substitution $y=\pi(n+a)^2x$ in each
summand proves CFD24 with every constant retained. This is the
shifted theta representation of the actual Hurwitz series before
multiplication by $s(s-1)/2$ or any passage to a cosine transform.

Its exact parameter derivative and its integral are
\begin{equation}\tag{CFD25}
\begin{gathered}
 \partial_tK_t(x)
     =-2\pi x\sum_{n=0}^{\infty}(n+a)e^{-\pi(n+a)^2x},\\
 \int_0^\infty \partial_tK_t(x)x^{s/2-1}\,dx
     =-s\pi^{-s/2}\Gamma(s/2)\zeta(s+1,a),
                     \qquad\operatorname{Re}s>0.
\end{gathered}
\end{equation}
Local uniformity gives the first derivative formula. For the second,
the sum of absolute integrals at $\sigma=\operatorname{Re}s>0$ is
\[
 2\pi\sum_{n\ge0}(n+a)
 \int_0^\infty x^{\sigma/2}e^{-\pi(n+a)^2x}\,dx
 =2\pi^{-\sigma/2}\Gamma(\sigma/2+1)
       \sum_{n\ge0}(n+a)^{-\sigma-1}<\infty.
\]
Integration term by term at complex $s$ and the recurrence
$\Gamma(s/2+1)=(s/2)\Gamma(s/2)$ prove the claimed identity.
In particular on $\operatorname{Re}s>1$ it is the derivative
of CFD24, and the convergent differentiated integral itself
continues the result to the larger stated half-plane.

The positive-real-$x$ integral has a relevant complex-parameter
restriction. The same formula is locally uniformly holomorphic
when
\begin{equation}\tag{CFD26}
 \operatorname{Re}a>|\operatorname{Im}a|.
\end{equation}
On each compact subset of this sector there is $c>0$ with
$\operatorname{Re}(n+a)^2\ge c(n+\operatorname{Re}a)^2$ for all
$n\ge0$. The preceding absolute-integral bounds therefore apply
with fixed additional constants. The gamma substitution with a
complex parameter follows from the holomorphic identity theorem
from positive real $a$; the principal logarithm satisfies
$\log((n+a)^2)=2\log(n+a)$ in this sector. This proves the
holomorphic extension there. On the larger half-plane
$\operatorname{Re}a>0$, the Hurwitz function has the continuation
proved in HTB4, but the unrotated kernel integral cannot be claimed
on that basis. For example $a=1+2i$ gives
$\operatorname{Re}(a^2)=-3$, and its first summand is
$e^{3\pi x}e^{-4\pi i x}$; the subsequent summands grow more
slowly or decay, so the positive-real-$x$ Mellin integral is not
absolutely convergent at infinity. This identifies the precise
analytic domain used by the direct kernel construction.

\subsection{The retained pair and its actual support domain}

On $\mathcal M_s\oplus\mathcal M_s$, the exact map
\begin{equation}\tag{CFD27}
 V(F,D)=(U(F+D),UD),\qquad
 V^{-1}(H,J)=(U^{-1}(H-J),U^{-1}J)
\end{equation}
is a complex-linear isomorphism by direct composition. It therefore
lifts to an isomorphism from
$G_L(\mathcal M_s\oplus\mathcal M_s)$ to
$G_L(\mathcal M_z\oplus\mathcal M_z)$, preserving its one common
label. For the actual common-top pair
$(F_t,\zeta-F_t)$ the image is $(U\zeta,U(\zeta-F_t))$.
Both summands are recoverable. Projection to the first entry
alone has linear kernel $\{(F,-F):F\in\mathcal M_s\}$ and
discards the deformation in this particular family.

This does not extend to the product of two split carriers with
independent labels by an unmodified triangular semimodule map.
There the first output label would be $\lambda\vee\mu$. For
any nonzero $D$ and two different labels $\lambda,\lambda'$,
the two inputs $((0,\lambda),(D,1_L))$ and
$((0,\lambda'),(D,1_L))$ have identical images
$((UD,1_L),(UD,1_L))$. Thus common-label CFD27 has its stated
inverse, while independent-label data require their original
labels as additional retained data. This is another explicit
domain distinction needed for a claim of faithfulness.

\clearpage

\appendix

% Complete source: THETA_HEAT_FULL_SUPPORT.tex
% Independent mathematical derivation. Prefix THS. Requires amsmath,amssymb,hyperref.
\section{The theta heat flow through the complete supported coefficient map}
\label{sec:theta-heat-full-support}

\paragraph{Scope correction, 23 September 2026.}
This chapter lifts the specified classical analytic heat kernel. It does
not derive that kernel from the entire absolute arithmetic base.
THS22 and THS25 explicitly show that its orbit and evaluated scalar
observable do not distinguish all support labels. The exact original
zeta and shifted-Hurwitz receiving formulas, retained boundary component,
exceptional modules, and inverse Mellin map are reconstructed in
ZTC1--21 and CFD1--27. Those comparisons, rather than the scalar
theta orbit alone, specify which original arithmetic object is retained.

\paragraph{Original source and actual reading.}
The original human source is Brad Rodgers and Terence Tao,
\href{https://arxiv.org/abs/1801.05914v5}{\emph{The De Bruijn--Newman
constant is non-negative}}, arXiv:1801.05914v5. The author file
\texttt{fmp-template.tex}, retained in
\texttt{zero\_inversion\_intake/original\_tex}, has SHA--256
\begin{center}\small
\texttt{25b015c0fb151f1613247d82cf53a6561a320d4dab2bb4f809482e7c477a5817}.
\end{center}
Lines 1--155 were read, including the exact definitions
\texttt{phidef}, \texttt{htdef}, \texttt{hoz} and \texttt{sas}.
This is a bounded reading, not a reading of the complete paper.
The arguments here use the displayed kernel and integral, and prove
all convergence, differentiation and coefficient statements below
directly. No threshold theorem about the zeros is invoked.
The complete coefficient map is ZIL11--12 of the included
\texttt{ZERO\_INVERSION\_FULL\_LATTICE\_REVIEW.tex}; its source at
the time of reading has SHA--256
\begin{center}\small
\texttt{eb87afa8b3ffac41ca0845bb1f2ba49d9dba14e6c6ac0db86d72d0e5963a50b9}.
\end{center}
Its lines 1--135 and 231--290 were read, including the complete
coefficient laws and the map with its retraction. The needed algebra
is proved again here. The existing included ZHR source was read in
full for its original coordinates; none of its threshold implications
is a premise of this calculation.

\subsection{The original kernel and a uniform bound for every moment}

Keep the original variables and all constants:
\begin{equation}\tag{THS1}
\begin{gathered}
 \Phi(u)=\sum_{n=1}^{\infty}\Phi_n(u),\qquad
 \Phi_n(u)=(2\pi^2n^4e^{9u}-3\pi n^2e^{5u})e^{-\pi n^2e^{4u}},
 \quad u\ge0,\\
 H_t(z)=\int_0^\infty e^{tu^2}\Phi(u)\cos(zu)\,du,
 \qquad (t,z)\in\mathbb C^2.
\end{gathered}
\end{equation}
Complex time extends the same integral, without changing its kernel.
Set
\begin{equation}\tag{THS2}
 a=e^{-3\pi},\qquad
 C_4=\frac{1+11a+11a^2+a^3}{(1-a)^5},\qquad C_\Phi=2\pi^2C_4.
\end{equation}
Then, for every $u\ge0$,
\begin{equation}\tag{THS3}
 0<\pi(2\pi-3)e^{9u-\pi e^{4u}}
       \le\Phi(u)\le C_\Phi e^{9u-\pi e^{4u}}.
\end{equation}
Indeed
$\Phi_n(u)=\pi n^2e^{5u}(2\pi n^2e^{4u}-3)e^{-\pi n^2e^{4u}}$
is positive, since $2\pi-3>0$. Its $n=1$ term is at least the
lower bound: $2\pi e^{4u}-3\ge(2\pi-3)e^{4u}$.
For the upper bound discard the negative term and use
$n^2-1\ge3(n-1)$ for $n\ge1$ and $e^{4u}\ge1$. Thus
\[
 \sum_{n\ge1}n^4e^{-\pi(n^2-1)e^{4u}}
 \le\sum_{n\ge1}n^4a^{n-1}=C_4.
\]
The last identity follows by applying $(a\,d/da)^4$ to
$\sum_{n\ge1}a^n=a/(1-a)$ and dividing by $a$; differentiation
is valid on $|a|<1$ by the geometric series. In particular the
kernel series converges absolutely, uniformly on every compact
$u$ interval, with the stated global bound.

For $A,B\ge0$ and every integer $p\ge0$, define the finite
explicit number
\begin{equation}\tag{THS4}
 \mathcal M_p(A,B)=\frac{C_\Phi}{2\pi}
 \exp\left(\frac{3A^2}{32\pi}
            +\frac{(B+9+p)^2}{8\pi}-\frac\pi2\right).
\end{equation}
The following pointwise and integral bounds hold:
\begin{equation}\tag{THS5}
\begin{gathered}
 u^p e^{Au^2+Bu}\Phi(u)
 \le C_\Phi\exp\left(\frac{3A^2}{32\pi}
                  +\frac{(B+9+p)^2}{8\pi}\right)
                e^{-\pi e^{4u}/2},\\
 \int_0^\infty u^p e^{Au^2+Bu}\Phi(u)\,du
                  \le\mathcal M_p(A,B).
\end{gathered}
\end{equation}
To prove this, use $u^p\le e^{pu}$, with $0^0=1$ in this
inequality, and the two Taylor bounds
$e^{4u}\ge(32/3)u^4$ and $e^{4u}\ge8u^2$.
Completing the square in $u^2$ and then in $u$ gives
\[
 Au^2-\frac\pi4e^{4u}
       \le Au^2-\frac{8\pi}{3}u^4\le\frac{3A^2}{32\pi},\qquad
 (B+9+p)u-\frac\pi4e^{4u}
       \le(B+9+p)u-2\pi u^2\le\frac{(B+9+p)^2}{8\pi}.
\]
The remaining half of the negative exponential proves the first
line. Finally $e^{4u}\ge1+4u$ gives
$\int_0^\infty e^{-\pi e^{4u}/2}\,du\le e^{-\pi/2}/(2\pi)$,
proving the second line. These estimates include every moment
order, every positive exponential-square weight, and any fixed
linear exponential weight.

\subsection{All differentiations and infinite heat sums}

For $\alpha,\beta\ge0$ integers and compact bounds
$\operatorname{Re}t\le A$ and $|\operatorname{Im}z|\le B$,
where $A,B\ge0$, one has
\begin{equation}\tag{THS6}
 \partial_t^\alpha\partial_z^\beta H_t(z)
  =\int_0^\infty u^{2\alpha+\beta}e^{tu^2}\Phi(u)
                      \cos(zu+\beta\pi/2)\,du,
 \qquad
 |\partial_t^\alpha\partial_z^\beta H_t(z)|
     \le\mathcal M_{2\alpha+\beta}(A,B).
\end{equation}
Indeed $|\cos w|,|\sin w|\le e^{|\operatorname{Im}w|}$.
Thus THS5 dominates every displayed integrand on those compact
sets. Difference quotients may be written as integrals of the
next derivative along a segment in a slightly larger compact
set; the same majorant applies. Dominated convergence therefore
proves the first complex derivative and its integral formula.
Iteration proves every mixed derivative and commutation of their
orders. Alternatively the exponential and cosine series about
any $(t_0,z_0)$ converge under the integrable majorant
$e^{(\operatorname{Re}t_0+R)u^2+(|\operatorname{Im}z_0|+S)u}
\Phi(u)$ on a bidisc of radii $R,S$. This also proves joint
entireness. All these arguments also allow the sum over $n$ in
THS1 to pass through the integral and derivatives: the sum of
the absolute integrands is bounded by the same expression
because every $\Phi_n(u)$ is positive. In particular
\begin{equation}\tag{THS7}
 \partial_tH_t(z)=-\partial_z^2H_t(z),\qquad
 H_t(-z)=H_t(z),\qquad
 H_{\bar t}(\bar z)=\overline{H_t(z)}.
\end{equation}
The last two identities follow directly from the integral and
the real positive kernel.

The exact Taylor expansion is
\begin{equation}\tag{THS8}
 H_t(z)=\sum_{j=0}^{\infty}h_j(t)z^{2j},\qquad
 h_j(t)=\frac{(-1)^j}{(2j)!}
          \int_0^\infty e^{tu^2}\Phi(u)u^{2j}\,du.
\end{equation}
For $|z|\le R$, the sum of the absolute integrands is
$e^{\operatorname{Re}t u^2}\Phi(u)\cosh(Ru)$, bounded by
THS5. This proves the expansion, absolute locally uniform
convergence, and the absence of every odd coefficient.
Every $h_j$ is entire in $t$ by THS6. For real $t$ there is the
explicit strict lower bound
\begin{equation}\tag{THS9}
 (-1)^jh_j(t)\ge
 \frac{\pi(2\pi-3)}{(2j)!(2j+1)}
          \exp\bigl(\min(t,0)-\pi e^4\bigr)>0.
\end{equation}
Restrict the integral to $0\le u\le1$. On that interval THS3
is at least $\pi(2\pi-3)e^{-\pi e^4}$ and
$e^{tu^2}\ge e^{\min(t,0)}$; integrate $u^{2j}$.
Thus every even Taylor coefficient is nonzero at every real
time. No such nonvanishing statement is assumed at complex time.

For arbitrary complex $t,r$, the coefficient heat identity is
\begin{equation}\tag{THS10}
 h_j(t+r)=\sum_{m=0}^{\infty}
       \frac{(-r)^m(2j+2m)!}{m!(2j)!}\,h_{j+m}(t).
\end{equation}
To verify absolute convergence before summing, THS8 and the
triangle inequality give
\begin{equation}\tag{THS11}
 \sum_{m\ge0}\frac{|r|^m(2j+2m)!}{m!(2j)!}|h_{j+m}(t)|
 \le\frac1{(2j)!}\int_0^\infty
      u^{2j}e^{(\operatorname{Re}t+|r|)u^2}\Phi(u)\,du<\infty.
\end{equation}
For a displayed bound replace the exponent coefficient by
$\max(\operatorname{Re}t+|r|,0)$ and apply THS5.
After interchanging sum and integral, the signs satisfy
$(-1)^m(-1)^{j+m}=(-1)^j$ and the remaining sum is
$\sum_m(ru^2)^m/m!=e^{ru^2}$. This proves THS10.
The same argument gives, for any finite number of successive
complex time steps $r_1,\ldots,r_d$, an absolutely convergent
multiple series with majorant
\begin{equation}\tag{THS12}
 \frac1{(2j)!}\int_0^\infty u^{2j}
 e^{(\operatorname{Re}t+\sum_{a=1}^d|r_a|)u^2}\Phi(u)\,du.
\end{equation}
Indeed the product of the absolute exponential series is
$\prod_a e^{|r_a|u^2}$. Finite partial products increase to
that positive bound; their integrals converge by THS5. Tonelli's
theorem followed by absolute convergence permits all orders of
summation. The binomial or multinomial identity then combines
the steps into $\sum_a r_a$, with their original signs.

There is also the function identity, with absolute locally
uniform convergence in every complex variable,
\begin{equation}\tag{THS13}
 H_{t+r}(z)=\sum_{m=0}^{\infty}\frac{(-r)^m}{m!}
                   \partial_z^{2m}H_t(z).
\end{equation}
Insert THS6:
\[
 \partial_z^{2m}\cos(zu)=(-1)^mu^{2m}\cos(zu).
\]
The absolute integral majorant is
\[
 e^{(\operatorname{Re}t+|r|)u^2+|\operatorname{Im}z|u}\Phi(u);
\]
THS5 applies uniformly on compact sets. This proves THS13,
and applies after every further $t,r,z$ derivative by increasing
the moment order. The group law on this actual family follows
either from THS12 or by evaluating its final integral.
In particular coefficient differentiation gives
\begin{equation}\tag{THS14}
 h_j'(t)=-(2j+2)(2j+1)h_{j+1}(t).
\end{equation}
For connection with the entire-function domain, every $a>0$
satisfies
\begin{equation}\tag{THS15}
 \sup_{z\in\mathbb C}|H_t(z)|e^{-a|z|^2}
 \le\int_0^\infty
        e^{(\operatorname{Re}t+1/(4a))u^2}\Phi(u)\,du<\infty.
\end{equation}
This follows from $|z|u\le a|z|^2+u^2/(4a)$ and
$|\cos(zu)|\le e^{|z|u}$. Thus the actual family belongs to
the space of entire functions satisfying every positive Gaussian
bound. This is a proved inclusion, without changing its kernel.

\subsection{Every original lattice label and the actual coefficient map}

Let $L$ be any bounded distributive lattice, with no completeness
assumption. Its bottom and top are $0_L,1_L$. Retain
\begin{equation}\tag{THS16}
\begin{gathered}
 S=G_L(\mathbb C)=\{(0,\lambda):\lambda\in L\}
                 \cup\{(c,1_L):c\in\mathbb C\},\\
 (c,\lambda)+(d,\nu)=(c+d,\lambda\vee\nu),\qquad
 (c,\lambda)(d,\nu)=(cd,\lambda\wedge\nu),\\
 \tau=(0,0_L),\quad e=(0,1_L),\quad
 1_S=(1,1_L),\quad\widehat c=(c,1_L).
\end{gathered}
\end{equation}
If an amplitude sum is nonzero at least one input label is top;
if a product is nonzero both input labels are top. Thus these
operations preserve the carrier. The ring laws and lattice
distributivity prove all semiring laws. No lower label is
identified with another. The construction also covers the
one-element lattice, when $\tau=e$ and $S$ is the original ring.

Keep the programme's tropical semifield with zero $K$ and put
$K_L=K\otimes_{\mathbb B}L$. Its zero and unit are
$\mathbf0_K,\mathbf1_K$. The map that sends nonzero elements
of $K$ to $1\in\mathbb B$ and zero to zero preserves addition
and multiplication: a tropical sum is zero precisely when all
its terms are zero and a product is zero precisely when a
factor is zero. Tensoring this map with $L$ proves the
well-defined retraction
\begin{equation}\tag{THS17}
\begin{gathered}
 \iota_L(\lambda)=\mathbf1_K\otimes\lambda,\qquad
 \rho_L\left(\sum_i b_i\otimes\lambda_i\right)
       =\bigvee_{b_i\ne\mathbf0_K}\lambda_i,\qquad
 \rho_L\iota_L=\operatorname{id}_L,\\
 \Phi_{\mathbb C}:S\longrightarrow\mathbb C\times K_L,
 \qquad\Phi_{\mathbb C}(c,\lambda)=(c,\iota_L(\lambda)).
\end{gathered}
\end{equation}
The retraction shows that $\iota_L$ is injective. Both coordinates
of $\Phi_{\mathbb C}$ preserve both operations, zero and unit,
and their equality determines both original coordinates.
Thus this is the injective original full-lattice map, not a
choice of one Boolean character.

For clarity about infinite sums, the following restricted
convergence is sufficient. A sequence $(c_m,\lambda_m)$ in $S$
has a sum here if $\sum_m|c_m|<\infty$ and the finite joins
$\lambda_0\vee\cdots\vee\lambda_M$ are constant for all
sufficiently large $M$. Its sum is the amplitude sum with that
eventual label. It belongs to $S$: if the amplitude sum is
nonzero at least one summand has nonzero amplitude and hence
top label. On these sums THS17 commutes with summation, because
the finite support sums in its target are exactly
$\iota_L(\lambda_0\vee\cdots\vee\lambda_M)$ and are likewise
eventually constant. This defines no arbitrary infinite join
in $L$ or $K_L$. Every infinite sum below satisfies this
restricted definition from its first summand.

The actual supported coefficient vector is
\begin{equation}\tag{THS18}
 \mathbf a(t)=(a_n(t))_{n\ge0},\qquad
 a_{2j}(t)=(h_j(t),1_L),\qquad a_{2j+1}(t)=\tau.
\end{equation}
At real times its even labels are forced by THS9 and the carrier
condition in THS16. The odd coefficients are structurally absent
from the cosine series and have the additive zero $\tau$.
For complex time the same labels are retained: if an even
coefficient happens to vanish it has value $e$, not $\tau$.
This gives a specified entire amplitude family with its unchanged
support vector, not an assertion of nonvanishing at complex times.

For every original supported time $x=(r,\lambda)\in S$, define
the upper triangular heat matrix on all degrees by
\begin{equation}\tag{THS19}
 U_x(i,j)=
 \begin{cases}
 1_S,&j=i,\\
 \displaystyle\left(\frac{(-r)^m j!}{m!i!},\lambda\right),
       &j=i+2m,\ m\ge1,\\
 \tau,&\text{otherwise}.
 \end{cases}
\end{equation}
Every entry lies in $S$. Indeed a nonzero off-diagonal amplitude
forces $r\ne0$, which forces $\lambda=1_L$. Equivalently that
entry is $\widehat{(-1)^m j!/(m!i!)}\,x^m$.
Thus if $r=0$, all off-diagonal amplitudes vanish but their
original label remains $\lambda$; they are not deleted.
For $x=(r,\lambda)$ and $y=(s,\mu)$, every matrix-product entry
is a finite sum and
\begin{equation}\tag{THS20}
 U_xU_y=U_{x+y},\qquad
 x+y=(r+s,\lambda\vee\mu),\qquad U_\tau=I.
\end{equation}
For an even gap $j-i=2m>0$, the amplitude of the product is
$\frac{j!}{i!}\sum_{b=0}^m(-r)^b(-s)^{m-b}/(b!(m-b)!)$
which is $(-r-s)^mj!/(m!i!)$. Its two endpoint terms have
labels $\lambda$ and $\mu$; each interior term has label
$\lambda\wedge\mu$. Their join is exactly $\lambda\vee\mu$,
including when some or every amplitude is zero. The diagonal
has its one unit term; an odd or negative gap has only absorbing
zero terms. This proves every entry and the identity statement.

For the actual vector, all coefficient sums converge in the
specified sense and satisfy the full supported heat law
\begin{equation}\tag{THS21}
 U_{(r,\lambda)}\mathbf a(t)=\mathbf a(t+r),
 \qquad t\in\mathbb C,\quad(r,\lambda)\in S.
\end{equation}
At row $2j$, its amplitude sum is THS10, absolutely convergent
by THS11. The $m=0$ summand already has label $1_L$, so every
finite partial sum has that same label. At every odd row every
source coefficient is $\tau$, and every product and partial sum
is $\tau$. This proves all amplitudes and every label.
For successive actions the absolute double or multiple sums
are THS12. Their finite support joins again have the displayed
labels because they include the diagonal contribution. Thus
THS20 acts on this vector with no change of summation order or
support completion left unproved.

In particular every zero-amplitude time, with its full original
label, acts by
\begin{equation}\tag{THS22}
 U_{(0,\lambda)}\mathbf a(t)=\mathbf a(t)
 \quad(\lambda\in L),\qquad
 U_e\mathbf a(t)=\mathbf a(t),\qquad U_\tau\mathbf a(t)=\mathbf a(t).
\end{equation}
These equalities do not equate the different operators on
arbitrary supported vectors. For example a vector whose only
nonzero entry is $a_2=1_S$ acquires coefficient $(0,\lambda)$
in degree zero under $U_{(0,\lambda)}$. On THS18 the even
degree-zero label, and every even tail label, is already top;
the same action adds no support. No label has been discarded.

Apply $\Phi_{\mathbb C}$ to every coefficient and to every
matrix entry; denote the resulting objects by
$\boldsymbol{\mathcal A}(t)$ and $\mathcal U_x$. Then
\begin{equation}\tag{THS23}
 \Phi_{\mathbb C}^{\mathbb N}(U_x\mathbf a(t))
   =\mathcal U_x\Phi_{\mathbb C}^{\mathbb N}(\mathbf a(t))
   =\boldsymbol{\mathcal A}(t+r),\qquad
 \mathcal A_{2j}(t)=(h_j(t),1_{K_L}),\quad
 \mathcal A_{2j+1}(t)=(0,0_{K_L}).
\end{equation}
Finite partial sums commute because THS17 is a semiring map;
their absolute amplitude limits and eventually constant support
limits commute by the explicit sum definition. Injectivity
retains every coefficient, including all supported zeros.

The coefficient heat differential also holds before and after
this map. On the fixed support vector in THS18, define the time
derivative by its proved entire amplitudes, retaining its even
top and odd bottom labels. The formal second derivative in $z$
has row $i$ equal to $\widehat{(i+2)(i+1)}\,a_{i+2}$.
Multiplication by $\widehat{-1}$ keeps the labels. THS14 proves
\begin{equation}\tag{THS24}
 \partial_t\mathbf a(t)=-D_z^2\mathbf a(t),\qquad
 \Phi_{\mathbb C}^{\mathbb N}(\partial_t\mathbf a(t))
      =-D_z^2\boldsymbol{\mathcal A}(t).
\end{equation}
This derivative is not an undeclared subtraction operation in
an arbitrary semiring. Each even coordinate lies in the exact
additive group $\mathbb C\times\{1_L\}$, whose identity is $e$;
each odd coordinate is the fixed singleton $\{\tau\}$.
Their product therefore has its coefficientwise complex vector
space structure with zero vector $(e,\tau,e,\tau,\ldots)$.
Differentiation there is ordinary coordinate differentiation;
the preceding integral bounds justify it for this family.
The inclusion into the full supported coefficient space retains
these identities and is not asserted to preserve its global
zero vector $(\tau,\tau,\ldots)$. This proves the precise meaning
of the differential in THS24 without changing any support.

Evaluation at any original supported point $y=(z,\nu)\in S$
also converges and gives
\begin{equation}\tag{THS25}
 \sum_{n\ge0}a_n(t)y^n=(H_t(z),1_L),\qquad
 \sum_{n\ge0}\Phi_{\mathbb C}(a_n(t))
                  \Phi_{\mathbb C}(y)^n=(H_t(z),1_{K_L}).
\end{equation}
Absolute convergence is THS8; the constant term is already
top-supported, so every partial support sum is top, including
$z=0$ and every allowed label $\nu$ of that zero. The two
equalities follow by THS17 and the ordinary Taylor expansion.
In particular an analytic zero has the original value $e$,
and its exact image $(0,1_{K_L})$, throughout this heat flow.

\begin{figure}[ht]
\centering
$\begin{array}{ccc}
 \mathbf a(t)&\xrightarrow{\ U_{(r,\lambda)}\ }&\mathbf a(t+r)\\[3pt]
 {\scriptstyle\Phi_{\mathbb C}^{\mathbb N}}\downarrow&&
             \downarrow{\scriptstyle\Phi_{\mathbb C}^{\mathbb N}}\\[3pt]
 \boldsymbol{\mathcal A}(t)&\xrightarrow{\ \mathcal U_{(r,\lambda)}\ }&
                         \boldsymbol{\mathcal A}(t+r)
 \end{array}$
\caption{The actual coefficient heat map, THS19--23. An even
coordinate retains label $1_L$ and its exact moment amplitude;
an odd coordinate is $\tau$. Every zero-amplitude time
$(0,\lambda)$ fixes this particular vector by THS22. All lattice
labels remain in the time object and the full coefficient map;
the diagram is not a replacement by a Boolean character.
The kernel and coordinates are Rodgers--Tao's THS1.}
\end{figure}

\subsection{The exact defect of every finite polynomial truncation}

Let $P_NH_t(z)=\sum_{j=0}^Nh_j(t)z^{2j}$ and let $P_N$ on
supported vectors replace all coefficients above degree $2N$
by $\tau$. At $0\le j\le N$ the amplitude of the discrepancy
between $P_NH_{t+r}$ and the polynomial heat evolution
$e^{-rD_z^2}P_NH_t$ is exactly
\begin{equation}\tag{THS26}
\begin{split}
 \Delta_{N,j}(t,r)
 &=\sum_{m=N-j+1}^{\infty}
       \frac{(-r)^m(2j+2m)!}{m!(2j)!}\,h_{j+m}(t)\\
 &=\frac{(-1)^j}{(2j)!}\int_0^\infty e^{tu^2}\Phi(u)u^{2j}
   \left(e^{ru^2}-\sum_{m=0}^{N-j}\frac{(ru^2)^m}{m!}\right)du.
\end{split}
\end{equation}
This follows by subtracting the finite beginning of the absolutely
convergent series THS10; the polynomial derivative formula gives
exactly those retained terms. For complex $t,r$, writing
$M=N-j+1\ge1$, the tail has the explicit estimate
\begin{equation}\tag{THS27}
 |\Delta_{N,j}(t,r)|
 \le\frac{|r|^M}{(2j)!M!}
      \int_0^\infty u^{2N+2}
          e^{(\operatorname{Re}t+|r|)u^2}\Phi(u)\,du.
\end{equation}
Use $\sum_{m=M}^{\infty}v^m/m!\le v^Me^v/M!$ for $v\ge0$:
the inequality $(M+\ell)!\ge M!\ell!$ proves it term by term.
The right side is finite and explicitly bounded by THS5.
If $t$ is real and $r>0$, then
\begin{equation}\tag{THS28}
             (-1)^j\Delta_{N,j}(t,r)>0
             \quad(0\le j\le N).
\end{equation}
The exponential tail in THS26 is strictly positive for every
$u>0$, and every other factor after removing $(-1)^j$ is
positive by THS3. Its finite integral is therefore positive.
This proves that finite polynomial truncation does not
intertwine the actual heat flow, even at the constant coefficient.
Injectivity of THS17 shows that the full coefficient map does
not remove this discrepancy. Both retained low coefficients
still have their top labels; their amplitudes differ by the
explicit nonzero quantity above.

Nevertheless the actual finite evolutions do converge to the
full flow, with the original coefficients. For $|z|\le R$ the
absolute double-series majorant for their terms is
\begin{equation}\tag{THS29}
 \int_0^\infty
 e^{(\operatorname{Re}t+|r|)u^2}\Phi(u)\cosh(Ru)\,du<\infty.
\end{equation}
Indeed the indices are $j,m\ge0$, and the absolute coefficient
integrands are bounded by
$u^{2j+2m}|r|^mR^{2j}/(m!(2j)!)$ times
$e^{\operatorname{Re}tu^2}\Phi(u)$. Summing gives exactly
the displayed majorant. The polynomial evolution retains the
triangle $j+m\le N$. These increasing triangles exhaust the
double series, so its absolute tails tend to zero, uniformly
on compact sets of $(t,r,z)$ by using their common THS5 bound.
Thus $e^{-rD_z^2}P_NH_t\to H_{t+r}$ locally uniformly, while
THS28 excludes exact intertwining at each finite $N$.
At each fixed supported coefficient its label is also eventually
the exact label in THS18. This states both the genuine limiting
map and the finite defect, with no assumed truncation identity.

\subsection{An explicit geometric bound for the entire discarded tail}

The convergence in THS29 has an explicit rate that retains the
original cutoff. Define
\[
 E_N(t,r,z)=H_{t+r}(z)-e^{-rD_z^2}P_NH_t(z).
\]
For every $\sigma>1$, $R\ge0$ and complex $t,r$, one has
\begin{equation}\tag{THS30}
 \sup_{|z|\le R}|E_N(t,r,z)|
 \le\sigma^{-(N+1)}
 \mathcal M_0\bigl(\max(\operatorname{Re}t+\sigma|r|,0),
                         \sqrt\sigma R\bigr).
\end{equation}
To prove this, the exact error is the absolutely convergent
double series of THS29 over $j+m>N$. For each such term,
$1\le\sigma^{j+m-(N+1)}$. Multiply its absolute integrand by
this bound, then sum over every $j,m\ge0$. The resulting upper
bound is precisely
\[
 \sigma^{-(N+1)}\int_0^\infty
 e^{(\operatorname{Re}t+\sigma|r|)u^2}\Phi(u)
                  \cosh(\sqrt\sigma R u)\,du.
\]
Since $\cosh v\le e^v$ for $v\ge0$, THS5 gives (THS30).
No coefficient, time parameter or cutoff has been rescaled in
the actual heat family; $\sigma$ is only an explicitly stated
parameter in the proved upper bound.

Every mixed derivative of this error has a corresponding exact
rate. For $\alpha,\beta,\gamma\ge0$, positive radii
$\delta_t,\delta_r,\delta_z$, complex $t_0,r_0$, and $R\ge0$,
put
\[
 A_*=\max\{\operatorname{Re}t_0+\delta_t
                 +\sigma(|r_0|+\delta_r),0\},\qquad
 B_*=\sqrt\sigma(R+\delta_z).
\]
Then
\begin{equation}\tag{THS31}
 \sup_{|z|\le R}
 |\partial_t^\alpha\partial_r^\beta\partial_z^\gamma
                  E_N(t_0,r_0,z)|
 \le
 \frac{\alpha!\beta!\gamma!}
      {\delta_t^\alpha\delta_r^\beta\delta_z^\gamma}
       \sigma^{-(N+1)}\mathcal M_0(A_*,B_*).
\end{equation}
The error is entire in all three variables by THS6, THS13
and its finite polynomial term. On the product of circles of
the three displayed radii about $(t_0,r_0,z)$, one has
$\operatorname{Re}t\le\operatorname{Re}t_0+\delta_t$,
$|r|\le|r_0|+\delta_r$ and $|z'|\le R+\delta_z$.
Apply (THS30) there, then apply the one-variable Cauchy
integral formula successively in the three variables. Each
formula contributes its exact factorial and inverse radius
power, proving (THS31). Its constants are finite by THS4 and
independent of $N$, so the bound tends geometrically to zero
for the original growing polynomial cutoff. The coefficient
supports at every retained even index are already top and
the odd indices are already $\tau$, as proved in THS18--29;
this amplitude error bound introduces no missing support limit.

The full supported heat law proved here retains the entire
amplitude family, all original lattice labels and their joint
map. It does not decide whether $H_0$ has nonreal zeros and
asserts no conclusion about RH.

\clearpage

% Complete source: HURWITZ_THETA_FLOW_BRIDGE.tex
% HTB: independent completed-Hurwitz / original-theta calculation.
% This file is a chapter source; the parent owns its cumulative wrapper.
\section{The exact completed Hurwitz--theta bridge over the full support lattice}
\label{sec:completed-hurwitz-theta-bridge}

\paragraph{Original-object correction, 23 September 2026.}
The auxiliary functions in this chapter are compared with the original
$\zeta$ and $F_t$; they are not replacements for them. HTB13 already
provides the fractional-module isomorphism at the exceptional points.
CFD1--10 retains its comparison with the original regular lattice and
all finite jets. ZTC13--16 proves the joint inverse with the boundary
component kept. Projection to $H_0$ alone discards the actual Hurwitz
deformation. The original arithmetic receiving formulas and both
Mellin endpoint terms are given in ZTC1--21.

\paragraph{Objects and source provenance.}
The shifted family below is the one in the supplied April 2026 programme
draft \emph{Split-Zero Support, Residue-Parity Reversal, Shifted-Hurwitz CUE
Defects, and Finite Phase-Space Enlargement}, printed byline ``Compiled
theorem draft'', source labels \texttt{thm:universal-flow} and
\texttt{thm:zero-motion}. The actual source is
\texttt{Chatnotes/globalization cue/}\allowbreak
\texttt{split\_zero\_cue\_shifted\_hurwitz\_phase\_space.tex},
lines 602--630 and 1003--1060; the exact retained copy is under
\texttt{hurwitz\_intake/programme\_source/}.
No human author is specified in those supplied
bytes, and none is inferred here. Its retained source has SHA-256
\texttt{52f2a6bccdb15213b1992abef9f8d6476d6e9ea5566ee7ed9d85c9a2d7cc2705}.
The heat family and every factor in its coordinate map are those of
Brad Rodgers and Terence Tao,
\href{https://arxiv.org/abs/1801.05914v5}{\emph{The de Bruijn--Newman
constant is non-negative}, arXiv:1801.05914v5}, original source
\texttt{fmp-template.tex}, equations \texttt{hoz}, \texttt{sas},
\texttt{phidef}, and \texttt{htdef}, lines 78--101. That source credits
N. C. de Bruijn's \emph{The roots of trigonometric integrals} (1950),
C. M. Newman's \emph{Fourier transforms with only real zeroes} (1976),
and the zero dynamics of G. Csordas, W. Smith, and R. S. Varga (1994).
Those earlier papers have not been independently read for this chapter;
their attribution is inherited explicitly from the original Rodgers--Tao
source. The theta identity used here is proved below, including its
Mellin factors. The full support carrier is the one in
\href{https://github.com/KokunoYumeto/zeta-function-research-reader/blob/a2851f9eb352e7e4376bc629de86fa4ed9c1c434/workbenches/splitzero-tandem/foundations/split-support-geometry-v11/split_support_geometry_arithmetic_curve_v11.tex}{\emph{Split Support Geometry}, v11,
Definition \texttt{def:lattice-split}}. Its operations and all maps needed
here are stated and verified in the final subsection.

\subsection{Two parameters and the original coordinate map}

Use $t$ for the Hurwitz shift and $h$ for heat time. They are independent
parameters. Put
\begin{equation}\tag{HTB1}
\begin{gathered}
 a=1+t,\qquad \operatorname{Re}a>0,\qquad
 F_t(s)=\zeta(s,a),\\
 A(s)=\frac{s(s-1)}2\pi^{-s/2}\Gamma(s/2),\qquad
 \Xi_t(s)=A(s)F_t(s),\\
 s(z)=\frac12+\frac{iz}{2},\qquad
 z(s)=-2i\left(s-\frac12\right),\qquad
 H_t^{\mathrm{Hur}}(z)=\frac18\Xi_t(s(z)).
\end{gathered}
\end{equation}
For complex $a$ in the indicated half-plane, $(n+a)^{-s}$ uses the
holomorphic logarithm with argument in $(-\pi/2,\pi/2)$. All subsequent
identities in $s$ are identities of meromorphic functions; values at
removable singularities mean their analytic extensions. Statements about
real parameters explicitly use $t>-1$.

Here is a direct continuation of the Hurwitz input, so that neither its
pole nor the endpoint values are assumed in the comparison. Define
Bernoulli polynomials by
\begin{equation}\tag{HTB2}
 \frac{x e^{b x}}{e^x-1}
 =\sum_{k=0}^{\infty}B_k(b)\frac{x^k}{k!},\qquad |x|<2\pi.
\end{equation}
For $\operatorname{Re}s>1$, termwise integration gives
\begin{equation}\tag{HTB3}
 \Gamma(s)F_t(s)
 =\int_0^{\infty}x^{s-1}\frac{e^{-a x}}{1-e^{-x}}\,dx.
\end{equation}
The interchange is absolutely convergent on compact subsets of
$\operatorname{Re}s>1$, $\operatorname{Re}a>0$: near zero the absolute
integrand is bounded by a constant times $x^{\operatorname{Re}s-2}$,
and at infinity it decays exponentially. Subtracting the first $N+1$
terms of HTB2 on $(0,1)$ gives, for $\operatorname{Re}s>-N$,
\begin{equation}\tag{HTB4}
\begin{split}
 F_t(s)=\frac1{\Gamma(s)}\bigg\{
 &\int_1^{\infty}x^{s-1}\frac{e^{-a x}}{1-e^{-x}}\,dx\\
 +&\int_0^1x^{s-1}
 \left(\frac{e^{-a x}}{1-e^{-x}}
       -\sum_{k=0}^{N}\frac{B_k(1-a)}{k!}x^{k-1}\right)dx\\
 +&\sum_{k=0}^{N}\frac{B_k(1-a)}{k!(s+k-1)}\bigg\}.
\end{split}
\end{equation}
The remainder is $O(x^N)$ locally uniformly in $a$. Differentiating in
$a$ and $s$ preserves locally integrable majorants on smaller compact
sets. Thus HTB4 gives joint meromorphic continuation, holomorphic in
$a$, with compatible values as $N$ increases.

The recurrence $\Gamma(s+1)=s\Gamma(s)$ and $\Gamma(1)=1$ show that
$\Gamma$ has residue $(-1)^n/n!$ at $-n$. Its reciprocal is entire with
simple zeros there and no other zeros: for example this follows directly
from the locally uniformly convergent Euler product
$1/\Gamma(s)=s e^{\gamma s}\prod_{n\ge1}(1+s/n)e^{-s/n}$, whose
nonvanishing away from the listed factors follows from the absolute
convergence of the logarithmic tail. In HTB4 these reciprocal zeros
cancel every displayed pole except the one at $s=1$. At $s=-n$ only
the term $k=n+1$ contributes to the value. Replacing $x$ by $-x$ in
HTB2 proves $B_k(1-a)=(-1)^k B_k(a)$. Consequently
\begin{equation}\tag{HTB5}
 \operatorname*{Res}_{s=1}F_t(s)=1,\qquad
 F_t(-n)=-\frac{B_{n+1}(a)}{n+1}\quad(n\ge0),\qquad
 \partial_t F_t(s)=-sF_t(s+1).
\end{equation}
For the last equality differentiate the convergent series on
$\operatorname{Re}s>1$; HTB4 and uniqueness of meromorphic continuation
extend it everywhere. In particular $F_t(0)=1/2-a$.

\subsection{The theta integral with all original constants}

For $x>0$, write
\[
 \psi(x)=\sum_{n=1}^{\infty}e^{-\pi n^2x},\qquad
 \Theta(x)=1+2\psi(x).
\]
The Gaussian identity gives
\begin{equation}\tag{HTB6}
 \Theta(x)=x^{-1/2}\Theta(1/x),\qquad
 \psi(1/x)=\sqrt{x}\,\psi(x)+\frac{\sqrt{x}-1}{2}.
\end{equation}
One direct proof of the required Gaussian identity is to periodize
$e^{-\pi x y^2}$. The periodized series and all its derivatives converge
uniformly on $[0,1]$. Its $k$th Fourier coefficient is
$\int_{\mathbb R}e^{-\pi x y^2}e^{-2\pi i k y}dy
=x^{-1/2}e^{-\pi k^2/x}$. This integral follows by scaling from
$\int e^{-\pi y^2}dy=1$ and differentiating its Fourier transform:
integration by parts gives $\widehat f'(\eta)=-2\pi\eta\widehat f(\eta)$,
with $\widehat f(0)=1$. The absolutely convergent Fourier series of the
periodization evaluated at zero proves HTB6.

Define the original kernel, without changing its argument or amplitude,
by
\begin{equation}\tag{HTB7}
\begin{split}
 \Phi(u)
 &=\sum_{n=1}^{\infty}
   \left(2\pi^2n^4e^{9u}-3\pi n^2e^{5u}\right)
       e^{-\pi n^2e^{4u}}\\
 &=e^u\left(2x^2\psi''(x)+3x\psi'(x)\right)_{x=e^{4u}}.
\end{split}
\end{equation}
If $g(u)=e^u\psi(e^{4u})$, direct differentiation and HTB6 give
\begin{equation}\tag{HTB8}
 \Phi(u)=\frac18\bigl(g''(u)-g(u)\bigr),\qquad
 g(-u)=g(u)+\sinh u,\qquad \Phi(-u)=\Phi(u).
\end{equation}
Indeed $g''-g=24e^{5u}\psi'(e^{4u})+16e^{9u}\psi''(e^{4u})$,
and $\partial_u^2-1$ annihilates $\sinh u$ and commutes with reflection.
For $u\ge0$, each summand in HTB7 is strictly positive because
$2\pi n^2 e^{4u}-3>0$. Moreover the series and each fixed number of its
derivatives are bounded by a constant times $e^{C u}e^{-\pi e^{4u}}$:
factor out $e^{-\pi e^{4u}}$ and bound the remaining sums of powers of
$n$ times $e^{-\pi(n^2-1)}$. Evenness gives the corresponding estimate
as $u\to-\infty$. These bounds justify all differentiations and
integrations that follow.

For $\operatorname{Re}s>1$, substitute $x=e^{4u}$ and integrate by parts
twice. The endpoint terms vanish, at infinity exponentially and at zero
by HTB6 and $\operatorname{Re}s>1$. The exact calculation is
\begin{equation}\tag{HTB9}
\begin{split}
 \int_{\mathbb R}\Phi(u)e^{(2s-1)u}du
 &=\frac14\int_0^{\infty}
       \left(2x^{s/2+1}\psi''(x)+3x^{s/2}\psi'(x)\right)dx\\
 &=\frac{s(s-1)}8\int_0^{\infty}\psi(x)x^{s/2-1}dx\\
 &=\frac{s(s-1)}8\pi^{-s/2}\Gamma(s/2)\zeta(s)
 =\frac14\xi(s).
\end{split}
\end{equation}
For clarity, with $b=s/2$ the integration-by-parts multiplier is
$\frac14(2b(b+1)-3b)=s(s-1)/8$. At zero HTB6 gives
$\psi(x)=\frac12x^{-1/2}-\frac12+O(x^{-1/2}e^{-\pi/x})$, and the
derivatives give the claimed vanishing boundary terms. The last
interchange in HTB9 is absolutely convergent for $\operatorname{Re}s>1$.
The left-hand side is entire by the preceding super-exponential bounds.
It therefore supplies the entire continuation of $\xi=A\zeta$ and,
by evenness, proves $\xi(s)=\xi(1-s)$. All its coefficients respect
complex conjugation because the kernel is real.

The original heat family is thus
\begin{equation}\tag{HTB10}
\begin{gathered}
 H_h^{\mathrm{heat}}(z)=\int_0^{\infty}e^{h u^2}\Phi(u)\cos(zu)\,du,
 \qquad (h,z)\in\mathbb C^2,\\
 H_0^{\mathrm{heat}}(z)=\frac18\xi(s(z))=H_0^{\mathrm{Hur}}(z),\qquad
 \partial_h H_h^{\mathrm{heat}}=-\partial_z^2 H_h^{\mathrm{heat}}.
\end{gathered}
\end{equation}
Joint holomorphy and the differential equation follow from locally
uniform integrable majorants for every $h,z$ derivative. To obtain the
middle identity, put $2s-1=iz$ in HTB9 and use evenness: the integral
over the whole real axis is twice the cosine integral. Hence the factor
is exactly $1/8$. Every $H_h^{\mathrm{heat}}$ is even; for real $h$ it
also satisfies $H_h^{\mathrm{heat}}(\bar z)
=\overline{H_h^{\mathrm{heat}}(z)}$.

\subsection{An invertible meromorphic map and its exact product}

Let $\mathcal M_s$ and $\mathcal M_z$ be the fields of meromorphic
functions on the entire $s$-plane and $z$-plane, respectively. The
following maps are everywhere defined as maps of these fields viewed
as complex vector spaces:
\begin{equation}\tag{HTB11}
 U:\mathcal M_s\longrightarrow\mathcal M_z,\qquad
 (Uf)(z)=\frac{A(s(z))}{8}f(s(z)),\qquad
 (U^{-1}G)(s)=\frac{8G(z(s))}{A(s)}.
\end{equation}
Since $A$ is a nonzero meromorphic function and $s,z$ are inverse affine
coordinates, these formulas prove complex linearity and both inverse
identities. With ordinary pointwise products $U$ is not multiplicative:
its two product expressions differ by a factor $A/8$. The exact
transported product and unit on the same meromorphic vector space are
\begin{equation}\tag{HTB12}
 (G\mathbin{\odot}J)(z)=\frac{8G(z)J(z)}{A(s(z))},\qquad
 \mathbf 1_{\odot}(z)=\frac{A(s(z))}{8},\qquad
 U(fq)=Uf\mathbin{\odot}Uq.
\end{equation}
All field axioms follow either by this inverse transport or by the
displayed multiplication formula. Thus $U$ is a unital field
isomorphism to $(\mathcal M_z,\odot)$, as well as a linear isomorphism
to the ordinary meromorphic receiver. These two assertions use
different, explicitly stated products. In particular the ordinary
constant function $1$ is not asserted to be its transported unit.

There is an exact local regularity rule. At a finite $s_0$, put
$z_0=z(s_0)$ and let $\nu_{s_0}(f)$ denote meromorphic order, with
$\nu(0)=+\infty$. Since the coordinate derivative $ds/dz=i/2$ is
nonzero,
\begin{equation}\tag{HTB13}
 \nu_{z_0}(Uf)=\nu_{s_0}(A)+\nu_{s_0}(f),\qquad
 \mathcal D_{s_0}
 =\{f:\nu_{s_0}(f)\ge-\nu_{s_0}(A)\}
 \xrightarrow[\;U\;]{\ \cong\ }\mathcal O_{z_0}.
\end{equation}
Here the domain is a fractional module of meromorphic germs and the
arrow is a module isomorphism over the affine coordinate isomorphism;
it is not claimed to be a subring with the ordinary product. At a point
where $A$ is a holomorphic unit, it restricts to an isomorphism of
holomorphic modules and satisfies the exact evaluation formula
\[
 \operatorname{ev}_{z_0}Uf
   =\frac{A(s_0)}8\operatorname{ev}_{s_0}f.
\]
At a zero or pole of $A$, HTB13, rather than an undefined product of
separate point values, specifies the evaluation domain.

For example, if $A$ is a unit at $s_0$ and $F_t$ has a zero there,
ordinary affine substitution, denoted $C(f)=f\circ s$, induces the
unital algebra isomorphism
\begin{equation}\tag{HTB14}
 \mathcal O_{s_0}/(F_t)
 \xrightarrow[\;C\;]{\ \cong\ }
 \mathcal O_{z_0}/(H_t^{\mathrm{Hur}}).
\end{equation}
Indeed $U F_t=(A\circ s)C(F_t)/8$ and $A\circ s$ is a unit, so the
two target principal ideals coincide. This proves equality of local
multiplicities and the correspondence of every jet, without changing
their coordinate scale. Explicitly, for $w=z-z_0$,
\[
 (Uf)(z_0+w)=\frac18\sum_{n\ge0}
       \left(\frac{i}{2}\right)^n w^n
       \sum_{j=0}^n\frac{A^{(j)}(s_0)f^{(n-j)}(s_0)}{j!(n-j)!}.
\]
Thus the complete finite jet map is triangular, with nonzero diagonal
$A(s_0)(i/2)^n/8$. It is $C$, not multiplication by that unit, which
gives the ordinary unital algebra arrow in HTB14.

\subsection{Both infinitesimal generators and their difference}

Define $D:\mathcal M_s\to\mathcal M_s$ by $Df(s)=s f(s+1)$.
All differential and shift operators in the next display have the
entire meromorphic field as their algebraic domain; no bounded-operator
or semigroup assertion on an unspecified normed space is being used.
Direct use of HTB11 gives
\begin{equation}\tag{HTB15}
\begin{split}
 (D_H G)(z):=(UDU^{-1}G)(z)
 &=c(s(z))G(z-2i),\\
 c(s)=\frac{sA(s)}{A(s+1)}
 &=\sqrt{\pi}\,\frac{s(s-1)}{s+1}
       \frac{\Gamma(s/2)}{\Gamma((s+1)/2)},\\
 \partial_tH_t^{\mathrm{Hur}}&=-D_HH_t^{\mathrm{Hur}}.
\end{split}
\end{equation}
The sign and displacement follow from $s(z-2i)=s(z)+1$, and the last
identity follows from HTB5. Gamma recurrence gives the removable values
$c(0)=-2$, $c(-1)=-2\pi$, and $c(1)=0$. The complete product in HTB15
is meromorphic: at negative odd integers below $-1$ the vanishing
multiplier may multiply a pole of the shifted function, and evaluation
is performed after multiplication. For instance, at $s=0$ the formula
will give $\partial_tH_t^{\mathrm{Hur}}(i)=1/8$, exactly as HTB20 below.

The original heat generator $B=-\partial_z^2$ has pullback
\begin{equation}\tag{HTB16}
 \mathcal L_H:=U^{-1}BU,\qquad
 \mathcal L_H f(s)=\frac1{4A(s)}\frac{d^2}{ds^2}\bigl(A(s)f(s)\bigr).
\end{equation}
This includes both powers from $ds/dz=i/2$ and the minus sign in $B$;
the original factors $8$ in $U$ and $U^{-1}$ cancel each other.
The actual pulled-back heat trajectory is
$f_h(s)=8H_h^{\mathrm{heat}}(z(s))/A(s)$, with $f_0=\zeta$ and
$\partial_hf_h=\mathcal L_H f_h$. No assertion is needed that an
operator exponential converges on every meromorphic function.

The two actual generators in the same $s$-space are $\mathcal L_H$ and
$-D$. Their exact difference is the operator
\begin{equation}\tag{HTB17}
\begin{split}
 \mathcal E f(s)
 &:=(\mathcal L_H+D)f(s)\\
 &=\frac14 f''(s)+\frac{b(s)}2f'(s)
   +\frac{b'(s)+b(s)^2}{4}f(s)+s f(s+1),\\
 b(s)&=\frac{A'(s)}{A(s)}
 =\frac1s+\frac1{s-1}-\frac{\log\pi}{2}
       +\frac12\frac{\Gamma'(s/2)}{\Gamma(s/2)},\\
 (B+D_H)U&=U\mathcal E.
\end{split}
\end{equation}
The expansion follows by applying the ordinary product rule twice to
HTB16. Every displayed coefficient is interpreted meromorphically;
apparent singularities from the expanded expression do not change the
domain just specified. HTB17 is the exact map of the generator defect,
not a declaration that the two flows are unrelated.

\subsection{All poles, the entire parameters, and reflection}

The only finite singularities of $A$ are simple poles at $s=-2m$,
$m\ge1$. The apparent pole at zero is removable and $A(0)=-1$;
$s=1$ is its only zero, and
$A(s)=(s-1)/2+O((s-1)^2)$ there. These statements follow from the
gamma pole calculation following HTB4, the nonvanishing of gamma
away from its poles, and the two polynomial factors of $A$. Hence
HTB5 gives
\begin{equation}\tag{HTB18}
\begin{gathered}
 \Xi_t(0)=t+\frac12,\qquad \Xi_t(1)=\frac12,\\
 a_m:=\operatorname*{Res}_{s=-2m}A(s)
   =\frac{2(-1)^m(2m+1)\pi^m}{(m-1)!},\\
 r_m(t):=\operatorname*{Res}_{s=-2m}\Xi_t(s)
   =\frac{2(-1)^{m+1}\pi^m}{(m-1)!}B_{2m+1}(1+t).
\end{gathered}
\end{equation}
These are all possible poles of $\Xi_t$. A listed point is a pole
exactly when $r_m(t)\ne0$; otherwise it is removable and its value is
$a_m\partial_sF_t(-2m)$. To check every constant, the residue of
$\Gamma(s/2)$ in the $s$-coordinate is $2(-1)^m/m!$ and
$(-2m)(-2m-1)/2=m(2m+1)$. Multiplying their product by
$F_t(-2m)=-B_{2m+1}(1+t)/(2m+1)$ gives HTB18. The endpoint at one
uses the residue $1$ of $F_t$, not its undefined point value.

In particular the first residue and the exact classification are
\begin{equation}\tag{HTB19}
\begin{gathered}
 r_1(t)=2\pi(1+t)\left(t+\frac12\right)t,\\
 t\in\mathbb R,\ t>-1:\qquad
 \Xi_t\text{ is entire}\quad\Longleftrightarrow\quad
 t\in\left\{0,-\frac12\right\},\\
 \Xi_0(s)=\xi(s),\qquad
 \Xi_{-1/2}(s)=(2^s-1)\xi(s).
\end{gathered}
\end{equation}
Indeed $B_3(a)=a(a-1/2)(a-1)$ follows by expanding HTB2. Thus the
first pole alone excludes every other $t>-1$. At zero, entireness was
proved in HTB9. At $t=-1/2$, the absolutely convergent identity
$\zeta(s,1/2)=2^s\sum_{n\ge0}(2n+1)^{-s}=(2^s-1)\zeta(s)$
extends meromorphically and proves entireness. This also proves that
no unexamined higher gamma pole invalidates the converse.

The same statements in the original receiver coordinate are
\begin{equation}\tag{HTB20}
\begin{gathered}
 z_m=(4m+1)i,\qquad
 \operatorname*{Res}_{z=z_m}H_t^{\mathrm{Hur}}(z)
 =\frac{r_m(t)}{4i}
 =\frac{i(-1)^m\pi^m}{2(m-1)!}B_{2m+1}(1+t),\\
 H_t^{\mathrm{Hur}}(i)=\frac{2t+1}{16},\qquad
 H_t^{\mathrm{Hur}}(-i)=\frac1{16}.
\end{gathered}
\end{equation}
The residue factor is $1/(8s'(z_m))=1/(4i)$; it is different from
the value factor $1/8$. In particular the tangent to the Hurwitz
trajectory at zero has a genuine pole:
\begin{equation}\tag{HTB21}
 \operatorname*{Res}_{s=-2}\left.\partial_t\Xi_t(s)\right|_{t=0}
 =\pi,\qquad
 \operatorname*{Res}_{z=5i}\left.\partial_tH_t^{\mathrm{Hur}}(z)
                                      \right|_{t=0}=\frac{\pi}{4i}.
\end{equation}
This is obtained by differentiating the explicit polynomial $r_1(t)$.
In contrast $\partial_h H_h^{\mathrm{heat}}|_{h=0}$ is entire by
HTB10. Thus HTB17 is nonzero even on the actual initial zeta function;
the difference is witnessed by a specified principal part, without a
change of coordinate, multiplication factor, or zero convention.

Define the meromorphic reflection defect by
\begin{equation}\tag{HTB22}
\begin{gathered}
 R_t(s)=\Xi_t(s)-\Xi_t(1-s),\qquad R_t(1-s)=-R_t(s),\\
 R_t(0)=t,\qquad R_t(1)=-t,\qquad
 H_t^{\mathrm{Hur}}(z)-H_t^{\mathrm{Hur}}(-z)=\frac18R_t(s(z)).
\end{gathered}
\end{equation}
For real $t>-1$, its identically vanishing locus is exactly $t=0$:
necessity follows by evaluation at zero, and sufficiency follows from
HTB9. At the other entire parameter,
$R_{-1/2}(s)=(2^s-2^{1-s})\xi(s)$, which is nonzero at $s=0$.
If $r_m(t)\ne0$, $R_t$ has simple poles at both $-2m$ and $1+2m$,
with residue $r_m(t)$ at each. The reflection changes the sign of a
residue and the subtraction changes it again; the two pole sets are
disjoint. For real $t$, $\Xi_t(\bar s)=\overline{\Xi_t(s)}$ and
therefore
$H_t^{\mathrm{Hur}}(-\bar z)=\overline{H_t^{\mathrm{Hur}}(z)}$.
It is not legitimate to replace the last identity by reality on the
real $z$-axis when the reflection defect is nonzero.

The reflection is also an exact operator on the source field:
\begin{equation}\tag{HTB23}
\begin{gathered}
 (J_A f)(s)=\frac{A(1-s)}{A(s)}f(1-s),\qquad J_A^2=I,
 \qquad UJ_A=J_zU,\quad (J_zG)(z)=G(-z),\\
 [D,J_A]f(s)
 =s\frac{A(-s)}{A(s+1)}f(-s)
 -(1-s)\frac{A(1-s)}{A(s)}f(2-s).
\end{gathered}
\end{equation}
Both assertions follow by substitution; their denominators are nonzero
elements of the meromorphic field. In particular
$R_t/8$ in the receiver is $(I-J_z)UF_t$, or $U(I-J_A)F_t$.
This gives a map of the entire reflection defect, including all its
principal parts. The heat generator commutes with $J_z$, whereas its
Hurwitz counterpart has exactly the commutator transported from
HTB23. The two projections $(I\pm J_z)/2$ retain the complete even
and odd meromorphic parts; no part is discarded in their direct sum.

\subsection{The exact first transverse velocity at a critical zero}

Let $\rho$ be a simple nontrivial zeta zero with
$\operatorname{Re}\rho=1/2$. This subsection concerns each such zero;
it assumes neither that every nontrivial zero is simple nor that every
nontrivial zero lies on the line. At $\rho$, $A$ is a holomorphic unit,
so HTB14 applies. Joint holomorphy from HTB4 and the implicit function
theorem give the unique local root $\rho(t)$ with
$F_t(\rho(t))=0$, $\rho(0)=\rho$. Differentiating that equation gives
\begin{equation}\tag{HTB24}
 \rho'(0)=c_1(\rho)
 =\frac{\rho\zeta(\rho+1)}{\zeta'(\rho)},\qquad
 z_\rho(t)=-2i\left(\rho(t)-\frac12\right),\qquad
 \operatorname{Im}z_\rho'(0)=-2\operatorname{Re}c_1(\rho).
\end{equation}
The numerator is nonzero: $\rho\ne0$ and the absolutely convergent
Euler product for $\zeta(\rho+1)$ has no zero because
$\operatorname{Re}(\rho+1)>1$. This proves $c_1(\rho)\ne0$, but does
not prove that its real part is nonzero.

At a zero $\xi'(\rho)=A(\rho)\zeta'(\rho)$. Reflection and conjugation
give $\xi'(1-\rho)=-\xi'(\rho)$ and
$\overline{\xi'(\rho)}=\xi'(\bar\rho)=-\xi'(\rho)$, so this derivative
is nonzero and purely imaginary. Similarly
$c_1(1-\rho)=\overline{c_1(\rho)}$. Differentiating HTB22 at the fixed
point $\rho$ therefore yields the exact identity
\begin{equation}\tag{HTB25}
\begin{split}
 \left.\partial_tR_t(\rho)\right|_{t=0}
 &=-\xi'(\rho)c_1(\rho)
     +\xi'(1-\rho)c_1(1-\rho)\\
 &=-2\xi'(\rho)\operatorname{Re}c_1(\rho),\\
 \left.\partial_t\bigl(H_t^{\mathrm{Hur}}(z_\rho)
                 -H_t^{\mathrm{Hur}}(-z_\rho)\bigr)\right|_{t=0}
 &=-\frac14\xi'(\rho)\operatorname{Re}c_1(\rho),
 \qquad z_\rho=z(\rho).
\end{split}
\end{equation}
Thus first-order reflection failure at this zero and its first-order
transverse velocity vanish together, with the original nonzero factor
shown. This proves an exact relation at every critical simple zero;
it does not extrapolate a finite numerical nonvanishing list to all
zeros. The meromorphic global defect has already been determined
without such an extrapolation in HTB18--HTB23.

\subsection{The retained arithmetic correction and its boundary kernel}

The full difference from the original arithmetic function is a
specified entire function before completion:
\begin{equation}\tag{HTB26}
\begin{gathered}
 D_t^{\mathrm{ar}}(s)=\zeta(s)-F_t(s),\qquad
 C_t(s)=A(s)D_t^{\mathrm{ar}}(s),\qquad
 J_t(z)=\frac18C_t(s(z)),\\
 F_t+D_t^{\mathrm{ar}}=\zeta,\qquad
 \Xi_t+C_t=\xi,\qquad
 H_t^{\mathrm{Hur}}+J_t=H_0^{\mathrm{heat}}.
\end{gathered}
\end{equation}
Both summands of $D_t^{\mathrm{ar}}$ have residue $1$ at their only
possible pole, so their difference is entire. There is a convergent
boundary-kernel formula on the larger half-plane
$\operatorname{Re}s>0$:
\begin{equation}\tag{HTB27}
 D_t^{\mathrm{ar}}(s)=\frac1{\Gamma(s)}
   \int_0^{\infty}x^{s-1}d_t(x)\,dx,
 \quad
 d_t(x)=\frac{e^{-x}(1-e^{-t x})}{1-e^{-x}},\quad d_t(0)=t.
\end{equation}
At infinity both exponential terms decay for $\operatorname{Re}(1+t)>0$;
at zero Taylor expansion gives $d_t(x)=t+O(x)$. These facts prove
absolute convergence and local uniformity on the indicated domain.
On $\operatorname{Re}s>1$ it is the difference of HTB3 for $a=1$ and
$a=1+t$, so it gives the stated continuation by uniqueness. All
remaining values are fixed by the entire continuation just proved.
Differentiation gives
$\partial_tD_t^{\mathrm{ar}}=sF_t(s+1)$, also directly from HTB5.

This correction retains every completed boundary contribution:
\begin{equation}\tag{HTB28}
\begin{gathered}
 C_t(0)=-t,\qquad C_t(1)=0,\qquad
 \operatorname*{Res}_{s=-2m}C_t=-r_m(t),\\
 C_t(s)-C_t(1-s)=-R_t(s),\qquad
 \partial_t C_t(s)=sA(s)F_t(s+1).
\end{gathered}
\end{equation}
Each statement follows by subtracting the corresponding functions in
HTB26, using entireness and reflection of $\xi$ proved in HTB9, or by
differentiation. The cancellation of all the principal parts and of
the entire reflection defect is an identity of the original analytic
functions. It does not assert a zero-location theorem for any summand.

\subsection{Every map over the original lattice and its faithful receiver}

Let $L$ be any nontrivial bounded distributive lattice. There is no
finiteness, Boolean-special-case, or one-point restriction here. For a
complex vector space $V$, define the full support semimodule
\begin{equation}\tag{HTB29}
\begin{gathered}
 G_L(V)=\{(0,\lambda):\lambda\in L\}\cup(V\times\{1_L\}),\\
 (v,\lambda)+(w,\mu)=(v+w,\lambda\vee\mu),\qquad
 (c,\alpha)(v,\lambda)=(cv,\alpha\wedge\lambda),\\
 \tau=(0,0_L),\qquad e=(0,1_L).
\end{gathered}
\end{equation}
The scalars are $G_L(\mathbb C)$ with its original addition and
multiplication. Closure follows because a nonzero amplitude forces
each participating amplitude that produces it to have top support.
The semimodule axioms follow componentwise from vector-space laws and
the two distributive laws of $L$. In particular an amplitude
cancellation has the join of the original labels, rather than being
replaced by $\tau$. If $V$ is a commutative algebra, its supported
multiplication is $(v,\lambda)(w,\mu)=(vw,\lambda\wedge\mu)$; it
has exactly the same distributivity verification.

For every complex-linear map $T:V\to W$ occurring above, define
\begin{equation}\tag{HTB30}
 T^{\uparrow}(v,\lambda)=(Tv,\lambda).
\end{equation}
This is well-defined because $v=0$ implies $Tv=0$. Linearity gives
preservation of addition and of the displayed scalar action, including
all zero amplitudes. Composition and inverse maps lift exactly. Thus
$U^{\uparrow}$ is a semimodule isomorphism, the generators and their
difference commute with it as in HTB15--HTB17, and both reflection
projections and their zero-amplitude fibres lift with their original
labels, as specified in HTB33. Their semimodule kernels, meaning the
inverse images of $\tau$, are the different fibres identified there.
Where the transported algebra of HTB12 is used, $U^{\uparrow}$ is
also a unital semiring isomorphism, with unit
$(A\circ s/8,1_L)$. Its induced map on prime spectra is the usual
preimage map: an ideal is prime exactly when its inverse image is
prime, since a bijective multiplicative map preserves and reflects
membership in a product. This applies to every support prime as well
as every arithmetic prime; no prime stratum is erased by the bridge.

To verify the original faithful comparison, let $K$ be the programme's
tropical semifield with zero and Boolean subsemiring $\mathbb B$.
Its map $\beta:K\to\mathbb B$ sends zero to zero and every nonzero
element to one. It preserves products, and preserves sums because a
sum in an idempotent semifield is zero only when both summands are
zero. Regard $L$ as the $\mathbb B$-algebra with addition $\vee$ and
multiplication $\wedge$, and set
\begin{equation}\tag{HTB31}
\begin{gathered}
 K_L=K\otimes_{\mathbb B}L,\qquad
 \iota_L(\lambda)=1_K\otimes\lambda,\qquad
 \rho_L=\beta\otimes\operatorname{id}_L,\qquad
 \rho_L\iota_L=\operatorname{id}_L,\\
 \Phi_V:G_L(V)\longrightarrow V\times K_L,\qquad
 \Phi_V(v,\lambda)=(v,\iota_L(\lambda)),\\
 (T\times\operatorname{id}_{K_L})\Phi_V
       =\Phi_W T^{\uparrow}.
\end{gathered}
\end{equation}
The tensor universal property proves the retraction formula. It proves
injectivity of $\iota_L$ and therefore of $\Phi_V$. The coordinate
laws verify addition and the $G_L(\mathbb C)$-action in the product
receiver. When $V$ is an algebra they also verify multiplication.
The last line is a coordinate equality, and supplies the commuting
square for every previously specified morphism, not just for a
scalar amplitude map.

For regular germs, evaluation has its own such square. At $s_0$ where
$A$ is a unit, HTB13 lifts to
\begin{equation}\tag{HTB32}
\begin{gathered}
 \operatorname{ev}_{z_0}^{\uparrow}U^{\uparrow}(f,\lambda)
 =\left(\frac{A(s_0)}8 f(s_0),\lambda\right),\\
 \Phi_{\mathbb C}\operatorname{ev}_{z_0}^{\uparrow}U^{\uparrow}
 =\left(\operatorname{ev}_{z_0}U\times\operatorname{id}_{K_L}\right)
       \Phi_{\mathcal O_{s_0}}.
\end{gathered}
\end{equation}
At zeros or poles of $A$ the first domain is replaced by the exact
fractional module $\mathcal D_{s_0}$ of HTB13; evaluation of $Uf$ is
defined there and its lifted square still holds. The unital local
algebra arrow of HTB14 likewise gives its full $G_L$ arrow, and
retains the support of every nilpotent jet or vanishing amplitude.

In particular, at a regular zero of an original top-supported analytic
function the evaluated element is $e=(0,1_L)$, whose image is
$(0,1_{K_L})$. It is not $\tau$, whose image is $(0,0_{K_L})$.
For a linear defect operator $T$, the entire inverse image of all
zero-amplitude labels is
\begin{equation}\tag{HTB33}
 \{x:T^{\uparrow}x\in\{(0,\lambda):\lambda\in L\}\}
      =G_L(\ker T),\qquad
 (T^{\uparrow})^{-1}(\tau)=\{\tau\},\qquad
 (T^{\uparrow})^{-1}(e)=\ker T\times\{1_L\}.
\end{equation}
The first equality is just $Tv=0$ with its unchanged label. The second
uses the defining fact that label $0_L$ permits only amplitude zero;
the third uses the definition of $e$. These are three different,
fully specified fibres.

Finally the correction identity HTB26 lifts without support loss:
\begin{equation}\tag{HTB34}
 (H_t^{\mathrm{Hur}},1_L)+(J_t,1_L)
   =(H_0^{\mathrm{heat}},1_L),\qquad
 \Phi(H_t^{\mathrm{Hur}},1_L)+\Phi(J_t,1_L)
   =\Phi(H_0^{\mathrm{heat}},1_L).
\end{equation}
This uses $1_L\vee1_L=1_L$ and its tensor image; it remains true at
points where the two amplitudes cancel, provided evaluations are
taken after the regular analytic sum. At an original zero, the sum
is supported zero. Separate summands with poles are not individually
evaluated. Thus the exact analytic correction, its reflection and pole
defects, the original theta flow, and the arbitrary-lattice receiver
form one commuting construction. The construction proves the stated
maps and defects; it supplies no conclusion that all original zeta
zeros lie on, or fail to lie on, the critical line.

\subsection{The exact boundary correction to the logarithmic receiver}

For every real $t\ge-1/2$, $H_t^{\mathrm{Hur}}(0)>0$. To prove this
without a numerical input, HTB7 and HTB10 give
$H_0^{\mathrm{heat}}(0)>0$. Since
$A(1/2)=-\Gamma(1/4)/(8\pi^{1/4})<0$, it follows that
$\zeta(1/2)<0$. At $a=1/2$ the identity used in HTB19 gives
$\zeta(1/2,1/2)=(\sqrt2-1)\zeta(1/2)<0$, while HTB5 gives
$\partial_a\zeta(1/2,a)=-\zeta(3/2,a)/2<0$ for every $a>0$;
the last sign follows from its absolutely convergent positive series.
Thus $F_t(1/2)<0$ for $a\ge1/2$, and the assertion follows from HTB1.

Fix such a $t$. On a sufficiently small disc about $z=0$, both
$H_0=H_0^{\mathrm{heat}}$ and $H_t=H_t^{\mathrm{Hur}}$ are
holomorphic and nonzero. Define there $c_t=J_t/H_0$, retaining the
correction $J_t$ of HTB26. Direct differentiation gives
\begin{equation}\tag{HTB35}
\begin{gathered}
 \frac{H_t'}{H_t}-\frac{H_0'}{H_0}
     =-\frac{c_t'}{1-c_t},\qquad
 W_t(z):=\frac{H_t(z)H_t(-z)}{H_0(z)^2}
     =(1-c_t(z))(1-c_t(-z)),\\
 b_n(t):=-\frac12[z^{2n-1}]
       \left(\frac{H_t'(z)}{H_t(z)}-
             \frac{H_t'(-z)}{H_t(-z)}\right),\quad n\ge1,\\
 b_n(t)=b_n(0)-\frac12[z^{2n-1}]\frac{W_t'(z)}{W_t(z)}.
\end{gathered}
\end{equation}
The brackets mean the indicated Taylor coefficient at zero. No branch
of a global logarithm is used. For the last identity differentiate
$H_t(z)H_t(-z)=H_0(z)^2W_t(z)$, using evenness of $H_0$. This proves
the correction to every coefficient of the even logarithmic receiver,
with no finite truncation. All terms are fixed by HTB27 and the
original theta integral.

The $b_n(t)$ are defined here as logarithmic coefficients, not as
uncorrected zero-only moments of a meromorphic function. The exact
distinction is its local divisor: if
$H_t(z)=(z-z_*)^{\nu}u(z)$ with $u(z_*)\ne0$, then
$H_t'/H_t=\nu/(z-z_*)+u'/u$. A zero of multiplicity $m$ contributes
$\nu=m$ and an actual pole contributes $\nu=-1$. HTB18--HTB20 specify
all those negative contributions and every cancellation. Thus any
global zero-only interpretation must retain exactly these pole
terms. HTB35 constructs the complete local correction to the
arithmetic receiver, and HTB31 applies to its linear coefficient
extractions with every support label unchanged. It does not assert
that the nonlinear logarithm is a linear morphism of split carriers.

\clearpage

% Complete source: HURWITZ_ALIGNMENT_FULL_MAP.tex
% Root derivation. HZA: full carrier, zero graph, and verified calculation scope.
\section{The aligned Hurwitz zero graph with its complete support map}
\label{sec:hurwitz-alignment-full-map}

\paragraph{Sources and actual objects.}
The original programme source is the retained April 2026 draft
\emph{Split-Zero Support, Residue-Parity Reversal, Shifted-Hurwitz CUE
Defects, and Finite Phase-Space Enlargement}, printed byline
``Compiled theorem draft'', source labels \texttt{thm:universal-flow},
\texttt{prop:non-euler}, and \texttt{thm:zero-motion}.
Its complete supplied TeX is retained with this calculation; its human
author is not specified and is not inferred here. The following proof
keeps its original $s,t$ coordinates. The operator construction in HPS
uses the arithmetic scaling of Jean-Beno\^it Bost and Alain Connes,
and the precise endomotive presentation of Alain Connes, Caterina
Consani and Matilde Marcolli,
\href{https://arxiv.org/abs/0806.2401v1}{\emph{Fun with $\mathbb F_1$}},
source equations \texttt{sigmatgenerators} and relations (c1)--(c4).
The full split carrier is the one in
\href{https://github.com/KokunoYumeto/zeta-function-research-reader/blob/a2851f9eb352e7e4376bc629de86fa4ed9c1c434/workbenches/splitzero-tandem/foundations/split-support-geometry-v11/split_support_geometry_arithmetic_curve_v11.tex}{\emph{Split Support Geometry}, v11, Definition \texttt{def:lattice-split}}.
The new claims below are proved here; numerical path labels have the
explicit, narrower meaning stated at the end.

\subsection{The full carrier and its finite-cardinality receiver}

Let $L$ be the original nontrivial bounded distributive lattice, and $R$
a commutative ring. Retain, without identifying any labelled zeros,
\begin{equation}\tag{HZA1}
\begin{gathered}
G_L(R)=\{(0,\lambda):\lambda\in L\}\cup(R\times\{1_L\}),\qquad
\tau=(0,0_L),\quad e=(0,1_L),\\
(a,\lambda)+(b,\mu)=(a+b,\lambda\vee\mu),\qquad
(a,\lambda)(b,\mu)=(ab,\lambda\wedge\mu),\\
K_L=K\otimes_{\mathbb B}L,\qquad
\iota_L(\lambda)=1_K\otimes\lambda,\qquad
\Phi_R(a,\lambda)=(a,\iota_L(\lambda)).
\end{gathered}
\end{equation}
Here $K$ is the programme's tropical semifield, and $\mathbb B$ its
Boolean subsemiring. The map $K\to\mathbb B$ sending zero to zero and
every nonzero element to one preserves addition and multiplication.
Tensoring it with $L$ gives $\rho_L:K_L\to L$ with
$\rho_L\iota_L=\operatorname{id}_L$. Thus $\Phi_R$ is injective,
preserves both operations and recovers every label, including labels
whose amplitude cancels. This supplies the actual comparison map,
rather than a count of extra points.

For finite $L$ and $n\ge1$, the exact underlying-set count is
\begin{equation}\tag{HZA2}
\left|G_L(\mathbb Z/n\mathbb Z)\right|=(n-1)+|L|,
\qquad
\sum_{n=1}^{\infty}\left|G_L(\mathbb Z/n\mathbb Z)\right|^{-s}
=\zeta(s,|L|),\quad \operatorname{Re}s>1.
\end{equation}
Indeed there are precisely $n-1$ nonzero amplitudes, each forced to
have top label, and precisely $|L|$ distinct zero-amplitude elements.
The sets are disjoint. Substituting $n-1+|L|$ in the absolutely
convergent sum proves the second equality. The quotient
$(a,\lambda)\mapsto a$ identifies all these zeros and gives the
ordinary count $n$. For $|L|=2$, the separated count gives
$\zeta(s,2)=\zeta(s)-1$. For arbitrary finite $L$ its shift is
$t=|L|-1$, with no change in the carrier (HZA1).

The count is only this specified receiver. It does not recover the
meet, join, or prime spectrum of $L$. For example, the four-element
chain and the four-element Boolean lattice have the same count;
they are not isomorphic because one is totally ordered and the other
has two incomparable elements. For infinite $L$ the displayed finite
cardinality series is not defined. The full map (HZA1) remains
defined. The real continuous shift below is an analytic family;
it is not an assertion that the cardinality of $L$ varies continuously.

\subsection{An analytic construction through the critical strip}

Put
\begin{equation}\tag{HZA3}
 F(s,t)=\zeta(s,1+t),\qquad
 h_t(x)=\frac{x e^{-(1+t)x}}{1-e^{-x}}.
\end{equation}
The branch of $\log(1+t)$ and of $\log(n+t)$ for complex $t$ is the
one on $\operatorname{Re}(1+t)>0$. On
$\operatorname{Re}s>0$, $s\ne1$, $\operatorname{Re}(1+t)>0$, the
following absolutely convergent expression constructs $F$:
\begin{equation}\tag{HZA4}
 F(s,t)=\frac1{\Gamma(s)}\left\{
 \frac1{s-1}+
 \int_0^1 x^{s-2}(h_t(x)-1)\,dx+
 \int_1^\infty x^{s-2}h_t(x)\,dx\right\}.
\end{equation}
Near zero, $h_t(x)=1+O(x)$, uniformly on compact parameter sets.
For large positive $x$, it is bounded by a constant times
$x e^{-c x}$, where $c>0$ is a lower bound for
$\operatorname{Re}(1+t)$ on the chosen compact set. The same bounds
after any fixed number of $s,t$ derivatives acquire only powers of
$x$ and $|\log x|$. These bounds prove absolute convergence and local
uniform differentiability under both integrals. For
$\operatorname{Re}s>1$, combine the two integrals with
$\int_0^1x^{s-2}\,dx=(s-1)^{-1}$ and expand
$(1-e^{-x})^{-1}=\sum_{j\ge0}e^{-jx}$. Absolute convergence gives
$\Gamma(s)\sum_{n\ge1}(n+t)^{-s}$. This proves (HZA4) in the original
half-plane and supplies its holomorphic continuation to the stated
domain. Its only possible pole there is $s=1$, with residue one.

Differentiating $h_t$ gives $\partial_t h_t=-x h_t$, so the two
integrals immediately give the exact flow equation
\begin{equation}\tag{HZA5}
 \partial_t F(s,t)=-s F(s+1,t),\qquad
 \partial_t^kF(s,t)=(-1)^k(s)_kF(s+k,t).
\end{equation}
The first equality follows from (HZA4) and the convergent Mellin
integral for $F(s+1,t)$; the second follows by induction, with
$(s)_k=s(s+1)\cdots(s+k-1)$ and $(s)_0=1$. Holomorphy in $|t|<1$
then gives the actual, locally uniformly convergent Taylor series
\begin{equation}\tag{HZA6}
 F(s,t)=\sum_{k=0}^{\infty}
       \frac{(-t)^k(s)_k}{k!}\zeta(s+k),\quad
       \operatorname{Re}s>0,\ s\ne1,\ |t|<1.
\end{equation}
In this domain the factors on the right introduce no additional pole:
for $k\ge1$ the argument has real part greater than one.
For any compact subset in $s$ and any smaller closed $t$ disc,
Cauchy's integral formula in a larger disc of radius less than one
proves the claimed uniform convergence. Thus no formal exponential
is being used as an unproved analytic operator assertion.

\subsection{Every nontrivial zero gives one exact local graph}

Let $\rho$ be any zero of $\zeta$ in $0<\operatorname{Re}\rho<1$,
and let $m$ be its order. This is a quantification over actual zeros;
no zero on or off a particular vertical line is supplied as an
assumption about all the zeros. The Euler product in the absolutely
convergent half-plane gives
\begin{equation}\tag{HZA7}
 B_\rho=\rho\zeta(\rho+1)\ne0,
 \qquad \partial_t F(\rho,0)=-B_\rho.
\end{equation}
For completeness, nonvanishing of that Euler product follows from
$\sum_p\sum_{k\ge1}|p^{-k(\rho+1)}|/k<\infty$: its exponential is
nonzero, and absolute expansion by unique prime factorization equals
the zeta series. Also $\rho\ne0$ in the stated strip.

There is a neighbourhood of $\rho$ and a unique holomorphic function
$\theta$ there, taking values near zero, for which the full local
zero set is
\begin{equation}\tag{HZA8}
 \{(s,t):F(s,t)=0\}=\{(s,\theta(s))\},\qquad \theta(\rho)=0.
\end{equation}
Here is a construction with the uniqueness included. Write
$B=\partial_t F(\rho,0)\ne0$. Choose sufficiently small closed
discs $|s-\rho|\le d$, $|t|\le r$ inside the holomorphy domain so
that $|1-B^{-1}\partial_tF(s,t)|\le1/4$ and
$|B^{-1}F(s,0)|\le r/2$. Such discs exist by continuity and
$F(\rho,0)=0$. The map $t\mapsto t-B^{-1}F(s,t)$ sends the
$t$ disc into itself and has Lipschitz constant at most $1/4$.
Starting from zero, its iterates converge uniformly to the unique
fixed point. Each iterate is holomorphic in $s$ on the open disc;
uniform convergence makes the limit holomorphic. A zero of $F$
in the product of discs is exactly a fixed point, proving (HZA8).

The order and leading coefficient of this graph are determined
without any simplicity assumption:
\begin{equation}\tag{HZA9}
 \theta(s)=(s-\rho)^m U_\rho(s),\qquad
 U_\rho(\rho)=\frac{\zeta^{(m)}(\rho)}{m!\rho\zeta(\rho+1)}\ne0.
\end{equation}
To prove this, the fundamental theorem of calculus gives
$F(s,t)=\zeta(s)+t A(s,t)$, where
$A(s,t)=\int_0^1\partial_tF(s,ut)\,du$ is holomorphic and
$A(\rho,0)=-B_\rho$. At the graph,
$\theta(s)=-\zeta(s)/A(s,\theta(s))$. Its denominator is a
holomorphic unit near $\rho$. Factoring the zero of $\zeta$ proves
both the order and the coefficient in (HZA9). Multiple zeta zeros
therefore give ramification of the projection to $t$, rather than a
failure to construct the total zero graph.

At a simple zero the graph can instead be written $s=\rho(t)$,
by the same contraction argument with the variables interchanged.
Implicit differentiation in (HZA5) proves
\begin{equation}\tag{HZA10}
 \rho'(0)=\frac{\rho\zeta(\rho+1)}{\zeta'(\rho)}.
\end{equation}
If this actual zero has $\operatorname{Re}\rho=1/2$ and the real
part of (HZA10) is nonzero, then
$\operatorname{Re}\rho(t)-1/2=t\operatorname{Re}\rho'(0)+O(t^2)$
for real $t$. Dividing by $t$ and using continuity proves a
transverse crossing for all sufficiently small nonzero $t$.
This conclusion is local. It places no restriction on a later
return of the same branch to the line.

\subsection{The supported graph and all its comparison maps}

For any open product of discs $D$ used above, let $R=\mathcal O(D)$.
The original function is the element $(F,1_L)$ of $G_L(R)$, and
its value at a point $x=(s,t)$ is obtained by the unital semiring
homomorphism
\begin{equation}\tag{HZA11}
\begin{array}{ccc}
 G_L(\mathcal O(D))&\xrightarrow{\Phi_{\mathcal O(D)}}&
                 \mathcal O(D)\times K_L\\
 {\scriptstyle(a,\lambda)\mapsto(a(x),\lambda)}\downarrow&&
       \downarrow{\scriptstyle(a,w)\mapsto(a(x),w)}\\
 G_L(\mathbb C)&\xrightarrow{\Phi_{\mathbb C}}&\mathbb C\times K_L.
\end{array}
\end{equation}
Evaluation preserves ring addition and multiplication, and the label
is unchanged, proving the left map preserves the displayed split
operations. Both composites are $(a(x),\iota_L\lambda)$, proving
commutativity on every element, including every labelled zero.
Consequently the zero graph just constructed is exactly
\begin{equation}\tag{HZA12}
 \{x:\operatorname{ev}_x(F,1_L)=e\}
 =\{x:\operatorname{ev}_x\Phi(F,1_L)=(0,1_{K_L})\}.
\end{equation}
Its value is supported zero, not $\tau$. All other zero sections
$(0,\lambda)$ remain present, evaluate to themselves, and are
recovered from the full map by $\rho_L$. The faithful comparison
does not impose an additional equation on $s$. The disc-algebra
prime operator in HPS gives the entire family another exact receiver;
HCR gives its residue-class receivers and their retained support
defects. These are calculations on the same $F$, with the displayed
maps, rather than a change of its zero convention.

\subsection{The actual bounded calculation and its limits}

The existing programme velocity certificate stores enclosures for
$V_j=2\operatorname{Re}\rho_j'(0)$. Its first 60 rows have 29
positive and 31 negative signs. The factor two must be retained when
comparing that file with (HZA10). The new script
\texttt{compute\_hurwitz\_paths.py} independently recomputes those
60 velocities at 45 decimal digits, checks each sign, and solves
for actual nearby Hurwitz zeros at $t=10^{-4}$ and $t=-10^{-4}$.
It checks residuals less than $10^{-33}$. These are numerical checks;
the pre-existing interval enclosures and the new non-interval
recomputations are recorded as different evidence.

The same script numerically follows the first five starting zeros
to $t=-1/2$ and $t=1$, with adaptive predictor/Newton steps and
80-digit endpoint refinement. At $t=1$, the first three endpoints
are, to the displayed decimal precision,
\[
 1.40778804065+23.32798775456i,\qquad
 0.41092175383+31.71541044219i,\qquad
 1.36768818112+38.41910556064i.
\]
The exact endpoint identities are
\begin{equation}\tag{HZA13}
 F(s,1)=\zeta(s)-1,\qquad
 F(s,-1/2)=(2^s-1)\zeta(s).
\end{equation}
The first follows by deleting the first term of the zeta series.
For the second, split that series into odd and even integers:
$\sum_{n\ge0}(n+1/2)^{-s}=2^s\sum_{n\text{ odd}}n^{-s}
=2^s(1-2^{-s})\zeta(s)$. Both identities extend from their
convergent half-plane by holomorphic continuation. The nonzero
factor $2^s$ is retained. The zeros of $2^s-1$ have
$s=2\pi i k/\log2$ for every integer $k$, including $k=0$.

The five paths in the figure are interpolations of numerical sample
points. They were initialized at known Riemann zeros; their common
value $\operatorname{Re}s=1/2$ at $t=0$ is therefore not an
independent proof for untracked zeros. The first and third numerical
paths reach the first two positive imaginary zeros of $2^s-1$;
the other three reach earlier Riemann zeros. No rigorous global
branch-identity certificate is claimed for this numerical path
tracking. Separate interval proofs in HXC certify two additional
transverse critical-line crossings at strictly positive shifts.
They show why the local conclusion (HZA10) must not be changed into
the assertion that a zero can cross only at an Euler-product parameter.

\clearpage
\end{document}
